% =====================================================================
%  TRIADE — édition académique v0.18.1 (manuscrit IMRAD, anglais)
%  L'édition v0.18.1 corrige les dix incohérences nommées par l'audit de
%  la vague 6 (instruction bqh4qlrqj) : plage des Entrées de clôture,
%  manifestations au slot, relation hint/RVA/adresse de crash, plage
%  exacte des ratios M1, note de non-additivité des médianes, décompte
%  des blocs crypto, réconciliation M4 (52.5/50.9, n=2), résidu PBT
%  (655/660), localisation de la table de checksums 0.1.1, table des
%  statuts complète (longtable). AUCUNE valeur mesurée n'est modifiée.
%  Corps du document : compilation Tectonic (XeTeX).
%  Les tableaux de données sont GÉNÉRÉS par les scripts d'extraction
%  (donnée unique : docs/JOURNAL_TRAVAIL.md + docs/DECISIONS.md + les
%  sorties brutes des instruments preuves_v110..v150 + les ledgers D204) et
%  \input ci-dessous — aucune donnée n'est tapée à la main ; les comptes du
%  journal vivent en macros (donnees/macros_v180.tex, régénérées à chaque
%  passe).
% =====================================================================
\documentclass[11pt,a4paper]{article}

\usepackage{fontspec}
\setmainfont{Liberation Serif}
\setsansfont{Liberation Sans}
\setmonofont{Liberation Mono}
\usepackage[french,english]{babel}

\usepackage{amsmath,amssymb}
\usepackage{geometry}
\geometry{a4paper, top=2.6cm, bottom=2.4cm, left=2.3cm, right=2.3cm}
\usepackage{graphicx}
\usepackage{xcolor}
\definecolor{liaison}{HTML}{3681A6}
\usepackage{booktabs}
\usepackage{tabularx}
\usepackage{longtable}
\newcolumntype{Y}{>{\raggedright\arraybackslash}X}
\usepackage{enumitem}
\setlist{noitemsep, topsep=2pt, leftmargin=1.4em}
\usepackage{microtype}
\usepackage{url}
\usepackage[numbers,sort&compress]{natbib}
\bibliographystyle{unsrtnat}
\setlength{\bibsep}{2pt}
\usepackage{fancyhdr}
\usepackage[colorlinks=true, linkcolor=liaison, citecolor=liaison,
            urlcolor=liaison, unicode=true,
            pdftitle={TRIADE: Enabling LLM Agents to Write Native Code by Hand Under Signed-Intent Governance (Academic Edition v0.18.1)},
            pdfauthor={Alexis Lopez}]{hyperref}
\pagestyle{fancy}
\fancyhf{}
\fancyhead[L]{\small\itshape TRIADE — academic edition v0.18.1}
\fancyhead[R]{\small\thepage}
\renewcommand{\headrulewidth}{0.3pt}

\newcommand{\code}[1]{\texttt{\small #1}}
\newcommand{\RQ}[1]{\textbf{RQ#1}}

% Les comptes EXACTS du journal (régénérés par extraire_donnees_v150.py —
% jamais un arrondi, jamais une divergence abstract/intro ; D190)
% Fichier GÉNÉRÉ (v0.18.0) — les comptes EXACTS des journaux récupérés.
\newcommand{\nbLignesJournal}{13{,}382}
\newcommand{\nbLignesWorklog}{11{,}856}
\newcommand{\nbLignesDecisions}{4{,}669}
\newcommand{\plageDecisions}{D1--D211}


\title{\textbf{TRIADE: Enabling LLM Agents to Write Native Code by Hand\\
Under Signed-Intent Governance\\
\large A longitudinal single-case study over eleven milestones, a
fourteen-unit two-arm field campaign, and two sealed applications}}
\author{Alexis Lopez\thanks{Independent researcher; corresponding author.
Project repository: \code{alexisfrlopez-debug/triade} (GitHub mirror).}}
\date{17 November 2026 — Academic edition v0.18.1 (the D204 campaign fully reconciled and decomposed; the RQ2 measurements unified under documented perimeters; the constant-time evidence reformulated as exact trace equality with the hardware runs reported; the unbounded proofs extended with a coverage lemma, a rule-by-rule TLA+ alignment and a 10{,}000-trace generative conformance campaign; the M4 verification time optimized without removing any control; the ten named inconsistencies of the v0.18.0 audit corrected --- the journal-closure range, the slot manifestations, the hint/RVA/crash-address relation, the exact M1-era ratio range, the median non-additivity note, the crypto-block count, the M4 reconciliation with the optimization's n{=}2 consigned, the PBT shrink residue, the checksum-table location, and the statuses table rendered complete)}

\begin{document}
\maketitle
\thispagestyle{fancy}

% ------------------------------------------------------------------
\begin{abstract}
\noindent Large language model (LLM) agents currently produce native code by
generating high-level source and delegating to a compiler toolchain; the
translation from intent to machine code passes through layers the agent
neither writes nor verifies. We report on TRIADE, a development protocol in
which an LLM agent writes native x86-64, RISC-V, and TLS~1.3 code
\emph{directly as assembly notation}, governed by signed layer-I
specifications, an encrypted and fragmented card that binds intentions to
byte ranges, and an epistemology that separates \emph{oracle} evidence
(differential emulation in the laboratory) from \emph{terrain} evidence
(execution on the developer's real machine). We present a longitudinal
single-case study of the protocol's validation programme: eleven milestones
(M0--M10, 26 October -- 11 November 2026), culminating in ``Loupe'', a real
Windows hex-editor application written entirely under the protocol and
delivered to its user --- and, closing the study, a pre-registered
two-arm field campaign (D204, 12--16 November 2026) in which the same
agent delivered every unit twice over a shared layer-I contract: once
as sealed hand notation (arm T) and once as compiled C (arm C) --- six
controlled blocks, then a second real application (``Passeur'', an
interception tool for the developer's authorized pentest) --- closed
green and signed at both arms. We quantify four research questions on data drawn
from a \nbLignesJournal-line append-only engineering journal (Entr'ees
consigned through 304; the campaign closures, Entr'ees 305--319,
carried in its \nbLignesWorklog-line companion task log and the campaign
ledger journal) and a \nbLignesDecisions-line append-only decision log
(\plageDecisions): (RQ1) the
escaped-defect rate of changes that had passed all laboratory controls ---
eleven of twelve field iterations surfaced at least one escaped defect,
each named to a class and a mechanism before any fix; the closing
campaign adds 22 witnessed deliveries in which zero escapes were
attributable to the sealed notation itself (six named kit-side escapes,
the unbound-leaf class recurring twice at the same slot, one pre-named
judge artifact); (RQ2) the measured
token cost per validated change --- re-measured end-to-end on the final
card with full distributions, the model's native tokenizer on every
instrumented call: medians of 1{,}277 tokens to locate a block, 3{,}374
to edit one, 3{,}674 to add one (the content-only v0.14.0 re-measures,
580/2{,}143/3{,}510, and the early single-shot M1 values, 51/719/1{,}088,
remain as explicitly labelled history --- the measurement perimeters
differ and are documented in the text) against a naive whole-disassembly
baseline of 300{,}304; the campaign's one fully instrumented
head-to-head unit reads 53{,}158 tokens for the notation arm against
67{,}475 for the compiled arm ($\times$1.27, the direction opposite to
the pre-registered expectation), the corrective spiral dominating both
arms (61.4\%/58.1\% of the totals);
(RQ3) the defect classes structurally invisible to emulation, which we
enumerate with their evidence (loader and PE-format semantics, ABI
shadow-space and SSE alignment, Windows error-code conventions,
harness-embedded assumptions) --- the campaign's recurrence being the
unbound import leaf that killed the compiled arm's field runs twice at
the same slot (the unbound IAT slot for memcpy at +0x151a0, whose
on-disk content --- the bare RVA of the hint/name-table entry, 0x152da
--- is the address the indirect call jumped to: the crash address is
the retained hint/name RVA itself), and the binding discipline that a
compiler normally provides silently; and (RQ4) the safety of hot-patching sealed
code --- zero post-patch regression across the milestone campaigns, with
signed version monotonicity and an armed anti-rollback check --- and a
closing campaign whose six field remediations moved zero bytes of the
sealed pieces (the gravures held across every iteration). The vault's four security properties (secret below threshold, rollback refusal, no destruction on wrong passphrase, blast radius bounded to one epoch) are now proved \emph{unbounded} over epochs by an inductive invariant preserved by every rule of the operational semantics, and the implementation is tested against the model by a 10{,}000-trace
generative conformance campaign (seed consigned, every failure path
covered, one named divergence class found, shrunk and analysed; zero gaps
elsewhere --- model-based conformance testing, renamed in the present
edition); a coverage lemma --- every reachable concrete state is
represented by one of the 288 abstract state shapes --- is stated and
discharged mechanically (SMT over unbounded integers, negative controls
included); each claim carries a systematic status --- proved unbounded,
proved bounded, measured, or awaiting a hardware run --- and no hardware
claim is made beyond what the developer's signed runs record. The study is
single-case (N=1), its oracle is derived from the same specifications it
checks, and the compiled-reference comparison shares the laboratory
harness; independent replication (a second developer, a second domain)
remains to be run. We state these limits explicitly. We claim as contributions the \emph{integration}
of contract-first notation, signed-manifest governance and oracle/terrain
verification into one pipeline for agent-first development --- not the
underlying cryptographic or testing primitives, which are standard.
\end{abstract}

\begin{otherlanguage}{french}
\begin{abstract}
\noindent Les agents de modèles de langage produisent aujourd'hui du code
natif en générant du source de haut niveau délégué à une chaîne de
compilation ; la traduction de l'intention vers le code machine traverse
des couches que l'agent n'écrit ni ne vérifie. Nous rapportons TRIADE, un
protocole de développement dans lequel un agent écrit du code natif
x86-64, RISC-V et TLS~1.3 \emph{directement en notation d'assemblage},
gouverné par des spécifications signées (couche I), une carte chiffrée et
fragmentée qui relie les intentions aux plages d'octets, et une
épistémologie qui sépare la preuve d'\emph{oracle} (émulation
différentielle au laboratoire) de la preuve de \emph{terrain} (exécution
sur la machine réelle du développeur). Nous présentons une étude de cas
unique longitudinale du programme de validation : onze jalons (M0--M10, du
26 octobre au 11 novembre 2026), culminant avec « Loupe », une véritable
application Windows (éditeur hexadécimal) écrite entièrement sous le
protocole et livrée à son utilisateur — et, pour clore l'étude, une
campagne terrain pré-enregistrée à deux bras (D204, 12--16 novembre 2026)
où le même agent a livré chaque unité deux fois sur un contrat de
couche I partagé : une fois en notation scellée à la main (bras T), une
fois en C compilé (bras C) — six blocs contrôlés puis une seconde
application réelle (« Passeur », un intercepteur pour le pentest
autorisé du développeur), close verte et signée aux deux bras. Nous quantifions quatre questions
de recherche sur des données tirées d'un journal d'ingénierie append-only
de \nbLignesJournal{} lignes (Entrées consignées jusqu'à 304 ; les clôtures
de campagne, Entrées 305--319, portées par le journal de tâches
compagnon de \nbLignesWorklog{} lignes et le journal de campagne) et d'un
journal de décisions append-only de \nbLignesDecisions{} lignes
(\plageDecisions) : (RQ1) le taux de défauts échappés des changements ayant
passé tous les contrôles de laboratoire — onze itérations de terrain sur
douze ont révélé au moins un défaut échappé, chacun nommé par classe et
mécanisme avant toute correction ; la campagne de clôture ajoute 22
livraisons témoignées dont zéro échappé imputable à la notation
scellée elle-même (six échappés côté kit nommés — la classe de la
feuille jamais liée récidivant deux fois au même slot —, un artefact
du juge pré-nommé) ; (RQ2) le coût mesuré en tokens par
changement validé — re-mesuré de bout en bout sur la carte finale avec
distributions, au tokenizer natif du modèle sur chaque appel instrumenté :
médianes de 1~277 tokens pour localiser, 3~374 pour éditer, 3~674 pour
ajouter (les re-mesures v0.14.0 au contenu seul, 580/2~143/3~510, et les
valeurs M1 51/719/1~088 demeurent historiques et étiquetées telles — les
périmètres de mesure diffèrent et sont documentés) contre une baseline
naïve de 300~304 ; l'unique unité à deux
bras entièrement instrumentée de la campagne lit 53~158 tokens pour le
bras notation contre 67~475 pour le bras compilé (×1,27, la direction
opposée à l'attente pré-enregistrée), la spirale corrective dominant les
deux bras (61,4 % / 58,1 % des totaux) ;
(RQ3) les classes de défauts structurellement invisibles à l'émulation,
énumérées avec leurs preuves (sémantique du chargeur et du format PE,
espace d'ombre et alignement SSE de l'ABI, conventions de codes d'erreur
Windows, hypothèses incrustées dans le harnais) ; et (RQ4) la sûreté du
patch à chaud de code scellé — zéro régression post-patch sur les
campagnes de jalons, avec monotonie de version signée et anti-rollback
armé. Les quatre propriétés de sûreté du coffre (secret sous le seuil, refus du retour en arrière, aucune destruction sur phrase erronée, rayon borné à une époque) sont désormais prouvées NON BORNÉES en époques par un invariant inductif préservé par chaque règle de la sémantique opérationnelle, et l'implémentation est testée contre le modèle par une campagne générative de conformité de 10~000 traces (graine consignée, tout chemin d'échec couvert, une classe d'écart unique trouvée, réduite et analysée ; zéro écart ailleurs — test de conformité model-based, renommé en la présente édition) ; un lemme de couverture — tout état concret atteignable est représenté par l'une des 288 formes d'état abstraites — est énoncé et déchargé mécaniquement (SMT sur entiers non bornés, contrôles négatifs compris) ; chaque affirmation porte un statut systématique — prouvé non borné, prouvé borné, mesuré, ou en attente de run matériel — et aucune revendication matérielle n'est faite au-delà de ce que les runs signés du développeur consignent. L'étude est un cas unique (N=1), son oracle dérive des mêmes
spécifications qu'il vérifie, et la comparaison au compilateur de référence
partage le harnais du laboratoire ; la réplication indépendante (un second
développeur, un second domaine) demeure à mener. Ces limites sont
énoncées explicitement. Nous
revendiquons comme contribution l'\emph{intégration} — notation à
contrats d'abord, gouvernance à manifestes signés, vérification
oracle/terrain — en un seul pipeline pour le développement piloté par
agent, et non les primitives cryptographiques ou de test sous-jacentes,
qui sont standards.
\end{abstract}
\end{otherlanguage}

% ------------------------------------------------------------------
\section{Introduction}\label{sec:intro}

An LLM-based agent that develops software today almost always works at the
level of high-level source code, and the mapping from that source to
machine code is delegated to a compiler toolchain. This division of labour
has an epistemological consequence: the agent's \emph{intent} is checked
against source text, while the \emph{bytes that execute} are produced by a
translation chain the agent does not write, does not see, and cannot
re-derive. Bugs in that mapping, and drift between what was specified and
what was produced, become a class of defects that no amount of source-level
review can exclude.

TRIADE (\emph{Triple Mirror: Intention, Bytes, Card}) is a protocol
designed for the opposite position: the agent writes native code
\emph{directly}, as assembly notation assembled by a deterministic,
verifiable assembler path --- there is no optimizing or opaque translation
chain; what the agent writes is, by construction, what executes
(assembly notation is verified by an exact round-trip against the emitted
bytes at every build). The protocol was developed and validated by a
single human developer directing an LLM agent, over eleven milestones
between 26 October and 11 November 2026, and is documented in a
\nbLignesJournal-line append-only engineering journal (Entr\'ees
consigned through 304; the campaign closures, Entr\'ees 305--319,
carried in its \nbLignesWorklog-line companion task log and the campaign
ledger journal at the archive), a \nbLignesDecisions-line append-only
decision log (\plageDecisions), and sealed artifacts whose
SHA-256 fingerprints are consigned at each sealing ceremony.

This paper is the first academic edition of that effort. Its purpose is
not to argue that hand-written machine code is superior to compiled code
in general; it is to report, with the data the protocol itself produced,
what happens when an LLM agent develops real native software this way ---
what it costs, what breaks, what laboratory verification could not see,
and what governance was needed to keep the process falsifiable. The
protocol's cryptographic and testing primitives (signature schemes,
secret sharing, differential testing, key ceremonies) are standard and
cited; the contribution claimed is their \emph{integration} into one
agent-first pipeline, together with the oracle/terrain epistemology and
the signed-intent governance described below.

The study is structured around four research questions:

\begin{itemize}
\item \RQ{1} \emph{Escaped-defect rate.} When a change passes every
laboratory control (differential oracles, harnesses at three fixed seeds,
structural checks, seal guards), how often does it still fail on the real
target machine, and at what granularity can the failure be named?
\item \RQ{2} \emph{Token cost per validated change.} What does it cost,
in LLM context tokens, to locate, edit, or add a block of sealed native
code under the protocol's card-based indexing --- measured end-to-end
on the final card, with distributions rather than single values?
\item \RQ{3} \emph{Emulation-invisible defect classes.} Which classes of
defects are structurally invisible to a Unicorn-based emulation harness,
and therefore escape to the terrain regardless of laboratory coverage?
Framed differently: a compiler silently guarantees several of these
classes --- RQ3 quantifies what must be re-derived, and learned, when the
compiler is removed from the agent's path.
\item \RQ{4} \emph{Hot-patch safety.} Can code be patched in place while
sealed artifacts remain byte-stable, with zero post-patch regression and
an armed rollback barrier, across ISAs?
\end{itemize}

The remainder of the paper describes the protocol and the study methods
(Section~\ref{sec:methods}), the results per research question
(Section~\ref{sec:results}), a discussion (Section~\ref{sec:discussion}),
threats to validity (Section~\ref{sec:threats}), related work
(Section~\ref{sec:related}), and declarations.

% ------------------------------------------------------------------
\section{Methods}\label{sec:methods}

\subsection{The TRIADE protocol in brief}\label{sec:protocole}

TRIADE organises development in three indissociable layers. The
\emph{intention layer} (layer I) holds one normalized specification per
code block: identifier, intent, pre/post-conditions and invariants, named
failure cases, verification vectors, fixtures, provenance and the
decisions that shaped it. The \emph{block layer} (layer B) holds the
assembly notation of each block, at fixed addresses, with deterministic
assembly and a mandatory round-trip check (assemble, disassemble,
re-assemble; any divergence refuses the build). The \emph{card layer}
(layer C) is an encrypted, fragmented, signed index (the \emph{card}) that
binds every block to its byte range, contract and fingerprints; only the
card is needed to locate or edit a block, never the whole binary.

Eleven rules (R1--R11) fix the discipline; the load-bearing ones for this
study are: R2 (\emph{written before code}: every block is specified before
any notation byte exists); R4 (\emph{stable contract}: a change may not
silently alter a block's contract); R5 (\emph{verify then write}: the
agent must re-derive the check before writing); R6 (crystallisation: an
artifact may only be sealed after ten consecutive identical rebuilds);
R7 (\emph{sealed inheritance}: inherited bytes live at their sealed
addresses, byte-identical); R10 (load-time verification: checksums per
block, signed version monotonicity, anti-rollback floor); R11 (no
optimization without measurement).

The development loop itself (update an intention, re-derive the oracle,
edit the notation, rebuild, verify, seal) and the seal ceremony model are
described in the protocol specification v0.13.0, which this edition
accompanies; we summarise only what the evaluation needs.

\paragraph{Governance: signed intent and append-only records.}
Every layer-I tree is condensed and signed with an Ed25519 key kept
strictly offline; a build refuses an unsigned layer I. Structural
decisions are consigned in an append-only decision log --- each entry has
an identifier, date, status, context, the retained choice, rejected
alternatives and amendment links; entries are never rewritten or
back-dated. The fourteen signed trees at the time of the study are verified
by a guard that must pass \emph{before} any build (14/14 at the last
run). Sealing a release is a \emph{fact of the card}: the sealed tag
lives in the card's build fragment; the artifact's bytes never move
(Table~\ref{tab:empreintes} lists the sealed artifacts and their
fingerprints).

\paragraph{The oracle/terrain epistemology.}
The protocol distinguishes two evidence regimes. \emph{Oracle evidence} is
produced in the laboratory: differential oracles derived from the layer-I
specifications (never exported, always regenerated from I), scenario
oracles that replay expected screens byte-for-byte, structural oracles
that read the composed image instruction by instruction, and Unicorn-based
emulation harnesses with Windows ABI shims. \emph{Terrain evidence} is
produced on the developer's real machine: the sealed artifact, a
scenario-loader script, and witnesses (byte-exact logs, showcase runs)
that are pushed back to an evidence vault and cross-checked (``recoup'')
against the laboratory's expectations. The study's unit of failure is the
\emph{named defect}: a failure is not counted until a class, a mechanism
and an evidence piece have been assigned to it, and every fix is a
documented surgical revision (kit-only or piece-only, with byte diffs).

\paragraph{The agent.}
The developer directed an LLM agent (Z.ai GLM family, ``the laboratory'')
through short natural-language instructions; the agent executed all
laboratory work (notation, oracles, harnesses, builds, ceremonies,
consignation) in a sandboxed Linux environment with no access to the
offline signing key. The complete transcript of the collaboration is the
append-only journal and work log themselves --- every session ends with a
consigned entry listing what was done, what was verified and what remains
--- plus the decision log; these, with the evidence vault, are the
primary sources of this paper. The agent is described here and in the
acknowledgements; it is not an author.

\subsection{The validation programme and the case study}\label{sec:cas}

The protocol was validated by eleven milestones, each with measured exit
criteria (Table~\ref{tab:jalons}); M0 established the artifact chain,
M2 hot patching on a real machine, M3 a million-case fuzzing campaign,
M4 an encrypted Schannel tunnel, M6/M7 portability to Linux and RISC-V,
M8 direct syscalls, M9 pure TLS~1.3 (28 blocks, official vectors from
FIPS~180-4, RFC~2104/4231, RFC~5869, RFC~7748, SP~800-38D, handshake
replayed on the RFC~8448 transcript, and a real encrypted session with a
real server), and M10 the full-scale test: \emph{Loupe}, a real Windows
hex-editor application --- 18 blocks in five families (file, digests,
rendering, navigation, editing) --- written entirely under the protocol,
delivered as a field kit to the developer's machine, and run twelve times
(Table~\ref{tab:runs}). Verification volumes per milestone are
consolidated in Table~\ref{tab:oracles}.

\subsection{Measurements}\label{sec:mesures}

\emph{Token budgets} were measured three times under two documented
perimeters. At M1 (decision D31/D34, the generation-1 card): the
single-shot content of each operation, cl100k\_base. At v0.14.0 (D195):
the content alone re-measured on the Loupe-era material (the plan entry,
the notation re-read and re-written, the specification and notation
written --- $n{=}10$ per operation), the agent \emph{reasoning} overhead
estimated at $\times 1.08$ from the era's surgical revisions (an
estimate, named as such, never confused with a measurement). At v0.17.0
(the pre-registered plan \code{ede70035\ldots}): the \emph{complete
instrumented round-trip} --- the real prompt (material verbatim plus
request prose) and the real response, counted with the model's native
tokenizer as read from the usage log of every call (the estimation
abandoned), cl100k\_base in parallel. The two re-measurements answer
different questions (what the material costs to read, against what the
full operation costs to live through); the perimeters are documented
wherever both values appear, and no direct cross-perimeter comparison
is drawn without the perimeter note. The M1 convention (D31/D34) was:
tokens to locate a block (read one \code{frag-loc} entry), to edit one
(additional \code{frag-con} + notation), to add one (additional
profile thresholds). \emph{Escaped defects} are counted per terrain
iteration under an operational rule defined once and applied uniformly:
a run counts toward the escaped-defect rate if at least one defect that
had passed every laboratory control in force at the time was named (a
class, a mechanism, an archived evidence piece) before any fix ---
kit-side and piece-side defects both count, the deliverable being the
kit+piece ensemble; runs~6 and~7 each surfaced two such defects and each
counts \emph{once} in the run rate (the rate measures escape incidence
per delivery, not per defect); the run-7 [D54-QUATRI\`EME] recurrence
is counted as a manifestation of the deposit-drift class first named in
the M4-era campaign, not as a new class; and the run-12 [T0-P\'ERIM\'E]
residue --- a stale banner string touching no gate, disclosed in the run
table --- is \emph{not} counted: it violates no control's contract and
alters no behaviour any control guards, and counting a text refresh as a
functional escape would misstate the rate's meaning.
\emph{Hot-patch safety} is measured as post-patch regressions (target:
zero) and as the byte-stability of sealed artifacts across patch
campaigns. All numbers reported in Section~\ref{sec:results} are
generated programmatically from the journal and decision log by the
extraction script shipped with this edition; the manuscript contains no
hand-typed data value.

% --- Tableaux générés depuis le journal (source unique de vérité) ---
% Fichier GÉNÉRÉ par extraire_donnees_v150.py — ne pas éditer à la main.
% Les cinq tableaux historiques : extraits/validés par le module extraire_donnees_v100.py (ancres du journal, empreintes mesurées).
% La table des preuves : PARSÉE depuis preuves_v110/v120/v130 — JAMAIS une troncature (la dette de la Table 6, soldée).

\begin{table}[htbp]\centering\footnotesize
\caption{Milestones M0--M10 of the TRIADE validation programme (26 Oct -- 11 Nov 2026). Every milestone was closed on measured exit criteria; sources: work journal entries and decision log.}
\label{tab:jalons}
\begin{tabularx}{\textwidth}{@{}ll lX@{}}
\toprule
\textbf{Milestone} & \textbf{Sealed tag} & \textbf{Date} & \textbf{Lab evidence (exit criteria met)} \\
\midrule
\textbf{M0} & \code{0.1.0} & 10-26 & round-trip 100\\%; executable < 4 KiB \\
\textbf{M1} & \code{0.2.0} & 10-26 & oracle 21/21; first crystallisation \\
\textbf{M2} & \code{0.3.0} & 10-27 & conformity 68/68; 8 instrumented iterations \\
\textbf{M3} & \code{0.4.0} & 10-28 & $10^6$ fuzzed cases $\times$2 seeds; 161 conformity controls \\
\textbf{M4} & \code{0.5.0} & 10-29 & 560-case SSPI scenario oracle $\times$4 seeds; 12 mutations caught \\
\textbf{M5} & \code{0.6.0} & 11-04 & per-build fingerprints active \\
\textbf{M6} & \code{0.7.0} & 11-05 & 7/7 conformity suite; layer-I diff zero \\
\textbf{M7} & \code{0.8.0} & 11-05 & round-trip 100 \% on new ISA; fence.i patch \\
\textbf{M8} & \code{0.9.0} & 11-05 & 61 harness greens; 131 socket-socle controls \\
\textbf{M9} & \code{1.0.0} & 11-07 & 28 blocks; 589 structural locks; handshake 54/54 $\times$3 seeds; cross-lock 86/86 \\
\textbf{M10} & \code{loupe 0.1.0 / 0.1.1} & 11-11 & 18 blocks; 159 controls $\times$3 seeds; cross-lock 26/26; seal guard 14/14 \\
\bottomrule
\end{tabularx}\end{table}

\begin{table}[htbp]\centering\footnotesize
\caption{Measured token budgets (M1 instrumentation, tokenizer cl100k\_base) against design targets, and generation-1 lineage volumes (D34). The lineage TOTAL counts specifications and delivered notations only --- drafts are excluded ($8{,}774 + 5{,}099 = 13{,}873$).}
\label{tab:budgets}
\begin{tabular}{@{}lrrl@{}}
\toprule
\textbf{Operation} & \textbf{Target} & \textbf{Measured} & \textbf{Source} \\
\midrule
Locate a block & 64 & 51 & D31/D34, M1 measurement (cl100k\_base) \\
Edit a block & 3000 & 719 & D31/D34, M1 measurement \\
Add a block & 6000 & 1,088 & D34, M1 (draft 1086, +2 drift re-measured) \\
Whole lineage (specs) & --- & 8,774 & D34, generation-1 card \\
Whole lineage (drafts) & --- & 3,973 & D34 \\
Whole lineage (delivered notations) & --- & 5,099 & D34 \\
Whole lineage (TOTAL) & --- & 13,873 & D34 \\
\bottomrule
\end{tabular}\end{table}

\begin{table}[htbp]\centering\footnotesize
\caption{Verification instruments and volumes per milestone. All numbers are drawn from the corresponding journal entries.}
\label{tab:oracles}
\begin{tabularx}{\textwidth}{@{}l X@{}}
\toprule
\textbf{Milestone} & \textbf{Instrument (measured volume)} \\
\midrule
\textbf{M1} & mini-libc oracle — 21/21 \emph{(D32; tag 0.2.0)} \\
\textbf{M2} & echo conformity + differential — 68/68; 8 iterations, 8 defect classes \emph{(D43; tag 0.3.0)} \\
\textbf{M3} & fuzz campaign + profile conformity suite — 1,000,000 cases $\times$2 seeds; 161 controls \emph{(D72/D73)} \\
\textbf{M4} & SSPI scenario oracle + per-block differential — 560 cases $\times$4 seeds; 10,000 fuzzed cases $\times$2 seeds per block; 165k leaf differentials \emph{(D68-D93)} \\
\textbf{M5} & build fingerprint registry + verifier — per-build sha256 registers \emph{(D134, D148)} \\
\textbf{M6} & profile conformity suite (Linux) — 7/7; lab recoup 22/22 \emph{(D154)} \\
\textbf{M7} & RISC-V round-trip + cross-profile lock — 100 \% round-trip; 29/29 recoup \emph{(D155)} \\
\textbf{M8} & sys harness on synthetic ntdll fixtures — 61 greens; 131 socle controls; 30/30 recoup \emph{(D160-D165)} \\
\textbf{M9} & tls13 structural oracle + handshake harness + cross-lock — 589 structural locks; handshake 54/54 $\times$3 seeds; cross-lock 86/86; 560+365+131 replayed \emph{(D166-D170)} \\
\textbf{M10} & Loupe: 3 Unicorn harnesses + structural oracle + seal guard — 159 controls $\times$3 seeds (61+53+45); 121 structural locks; cross-lock 26/26; seal guard 14/14; crystallisation 26 x 10/10 \emph{(D172-D181)} \\
\bottomrule
\end{tabularx}\end{table}

\begin{table}[htbp]\centering\scriptsize
\caption{The twelve terrain runs of the Loupe field test (M10). Every failure was \emph{named} (a class, a mechanism, an evidence piece) before being fixed; every fix was documented as a surgical revision. The sealed artifacts never changed after their sealing run. Sources: journal entries 204--222, decisions D178/D181.}
\label{tab:runs}
\begin{tabularx}{\textwidth}{@{}r l X X@{}}
\toprule
\textbf{\#} & \textbf{Kit} & \textbf{Outcome (named defect)} & \textbf{Documented fix} \\
\midrule
1 & \code{dev.1} & loader refusal 0xc000007b. Missing BASERELOC: ASLR image without relocation table (lab harness parses, never loads) & dev.2: minimal 12-byte BASERELOC, 7 bytes in 6 ranges \\
2 & \code{dev.2} & same refusal 0xc000007b. Necessary-not-sufficient: page hole 0x3000-0x4000 + massive BSS outside the proven envelope of 8 prior artifacts & dev.3: image within envelope (VS==RS, no BSS) \\
3 & \code{dev.3} & loader mapped the image (decisive step); T0 witness failed. [CADRE-NE]: .data initializer wrong (sentinel at +8 - a SIZE read as an offset instead of +528) & dev.4: initializer = exact mirror of harness CADRE\_VIERGE; guard healed (control derives from plan, never from the piece it guards) \\
4 & \code{dev.4} & T0 all green; T1 named -6; showcase died instantly. Stack-argument class: msx64 args 5+ at home space instead of [rsp+0x20]+ (loupe-01/05/06) + SSE alignment (loupe-10/11) & dev.5: slots fixed; msx64 ABI audit born at build (static, catches SSE class); falsification control V7 added \\
5 & \code{dev.5} & first run acquitting the piece; T1 named -7 VERROUILLE. [SESSION]: kit itself - scenario-as-loader holds loupe.exe as a live section, Windows refuses gle 32 to open it for write & dev.6 KIT-ONLY (artifact unchanged): byte-copy shadow temoin-session.bin \\
6 & \code{dev.6} & two defects from one witness. False save error + missing showcase thread; [HINSTANCE-PEB]: hInstance read from PEB instead of GetModuleHandleW(NULL) & piece dev.6 (r12 = GetModuleHandleW(NULL)) + kit dev.7 \\
7 & \code{dev.7} & T0-T5 all green (save proven on terrain); then two defects. [EAX-SANS-SIGNE] at kit (unsigned 32-bit reads of negative codes) + [D54-QUATRIEME] at piece (4 contiguous FirstThunk ranges without NULL separators) & piece dev.7 (IAT relocated 0x3100, 4 ranges separated by NULL, 63 call sites repointed) + kit dev.8 \\
8 & \code{dev.8} & T0-T7 all green for the FIRST time (functional contract fully proven). Probe died on its own Python; showcase still mute (5th death) & kit dev.9 (probe made load-bearing) \\
9 & \code{dev.9} & T0-T7 all green second time; 6th showcase death named by system then by lab. [OFT-QWORD]: 4 OFT tables at QWORD stride instead of DWORD - the loader walks 10.18.6.3 entries at DWORD stride; first cause of ALL deaths found on the piece bytes & piece dev.8 (89 bytes in 52 ranges, sealed 0.1.0 candidate) + kit dev.10 \\
10 & \code{dev.11} & T0-T7 all green third time; SHOWCASE ALIVE FIRST TIME (0 and 0, no WER). [FREE-SANS-TYPE] at kit: FreeLibrary called without argtypes, 64-bit HMODULE truncated to C int - ArgumentError killed cascade 205 & kit dev.12 (argtypes=[c\_void\_p], except Exception widened, probe armored try/finally) \\
11 & \code{dev.12} & T0-T7 all green 4th time; cascade 205 complete 1st time; showcase alive 2nd time; green holds entire. [SONDE-EXE] named by lab on pieces: structural false negative - library loader NEVER resolves an EXE's imports; liaison proven by the living showcase & seal loupe 0.1.0 on recouped run 39/39 (D178): notation debt paid (587$\rightarrow$0), crystallisation 26x10/10, anti-rollback armed \\
12 & \code{dev.1} & T0-T7 all green 5th time; cascade 205 2nd time; showcase alive 3rd time (0 and 0, first of piece 0.1.1). [T0-PERIME] residual named (dev.12-era banner text - no gate touched) & seal loupe 0.1.1 (D181): [R10-PERIMEE] debt paid at birth (17 sums refreshed), card generation 17, artifact intact \\
\bottomrule
\end{tabularx}\end{table}

\begin{table}[htbp]\centering\footnotesize
\caption{Sealed artifacts (SHA-256, verified byte for byte against the card and journal at edition time; the campaign pieces and instruments added at v0.17.0, each re-verified on the bytes at edition time).}
\label{tab:empreintes}
\begin{tabular}{@{}l r l@{}}
\toprule
\textbf{Artifact} & \textbf{Bytes} & \textbf{SHA-256 (first 16 hex)} \\
\midrule
\code{passeur.exe v0.2 (arm T)} & 57,856 & \code{4a7229e7418e8276\ldots} \\
\code{passeur-c.dll v0.2 (arm C)} & 97,937 & \code{fe4ddfc9842fb29b\ldots} \\
\code{pilot EXE (installer)} & 99,285,377 & \code{34b9764131e3b3a5\ldots} \\
\code{pilot source v2} & 216,080 & \code{40d23f03fdaaf116\ldots} \\
\code{pilot post (all runs)} & 67,514 & \code{50f7bc86dcca67cc\ldots} \\
\code{pilot post (campaign)} & 55,357 & \code{c2dede15894f4eeb\ldots} \\
\code{loupe.exe 0.1.1} & 127,488 & \code{c111510535767a40\ldots} \\
\code{loupe.exe 0.1.0} & 127,488 & \code{054de960c31425d0\ldots} \\
\code{proxy.exe 1.0.0} & 65,024 & \code{446279f3ab356900\ldots} \\
\bottomrule
\end{tabular}\end{table}

\begingroup\scriptsize
\begin{longtable}{@{}p{2.4cm} p{2.7cm} p{5.4cm} p{3.6cm}@{}}
\caption{Strengthened evidence across editions, every row generated from the raw instrument outputs (no value typed by hand): the six v0.11.0 instruments (D185, consigned seeds), the v0.12.0 formal model and seal analysis (D187, D189), the v0.13.0 corrections (the audited benchmark, the corrected seal statistics), the v0.14.0 empirical programme, and the v0.15.0 unbounded proofs and refinement.}\label{tab:preuves}\\
\toprule
\textbf{Instrument} & \textbf{Target} & \textbf{Measured} & \textbf{Raw source} \\
\midrule
\endhead
\bottomrule
\endfoot
Mutation testing (16 mutants/block, seed 20261111) & loupe domain (18 sealed blocks) & 230/288 killed (79.9\%); 58 survivors named & \code{preuves\_ \allowbreak{}v110/\allowbreak{}mutation.csv} \\
\addlinespace[1.5pt]
Mutation testing (16 mutants/block, seed 20261111) & tls13 domain (28 sealed blocks) & 397/448 killed (88.6\%); 51 survivors named & \code{preuves\_ \allowbreak{}v110/\allowbreak{}mutation.csv} \\
\addlinespace[1.5pt]
Mutation testing (16 mutants/block, seed 20261111) & sys domain (6 sealed blocks) & 67/96 killed (69.8\%); 29 survivors named & \code{preuves\_ \allowbreak{}v110/\allowbreak{}mutation.csv} \\
\addlinespace[1.5pt]
Mutation testing (same campaign) & cryptographic families (sha-256, hmac/hkdf, x25519, aes-gcm, md5) & 100\% killed (224/224 mutants, 0 survivors) & \code{preuves\_ \allowbreak{}v110/\allowbreak{}mutation.csv} \\
\addlinespace[1.5pt]
Coverage-guided fuzzing (seed 20261111) & loupe-01 (untrusted input path) & 72/119 instructions covered (60.5\%); 0 crashes; 0 canary violations; corpus 18 & \code{preuves\_ \allowbreak{}v110/\allowbreak{}fuzzage.json} \\
\addlinespace[1.5pt]
Coverage-guided fuzzing (seed 20261112) & http-10 (input parser, sealed 0.4.0) & 358 instructions covered; 0 crashes; 0 named faults; corpus 17 & \code{preuves\_ \allowbreak{}v110/\allowbreak{}fuzzage.json} \\
\addlinespace[1.5pt]
Bounded static memory safety (abstract interpretation) & 52 sealed blocks & 2\,098 accesses classified; 1\,181 stack accesses verified, 0 out-of-frame; 263 contract-mediated; 191 constant [rip]; 463 unresolved (named) & \code{preuves\_ \allowbreak{}v110/\allowbreak{}surete\_ \allowbreak{}memoire.csv} \\
\addlinespace[1.5pt]
Constant-time (emulation, exact trace equality; N=60/class, grain 20261111) & tls13-01..04 (sha-256), 05 (hmac), 06 (hkdf), 13/14 (x25519), 15..19 (aes-gcm), loupe-06 (md5) & ONE instruction count per target across all 120 measurements: sha-256 4\,469; hmac 18\,371; hkdf 18\,407; x25519 exactly 1\,238\,959 for every key; aes-gcm 11\,695; md5 132\,509 --- traces bit-identical, Welch $t$ undefined (the v0.11.0 ``$|t|{=}0$'' withdrawn; the hmac/aes-gcm counters had never run) & \code{preuves\_ \allowbreak{}v110/\allowbreak{}temps\_ \allowbreak{}constant.csv} + the v0.18.0 replay \\
\addlinespace[1.5pt]
Oracle independence (fresh seed 20261113) & MD5/SHA-256 tables re-derived (decimal arithmetic, cube roots); RFC 1321 gold vectors spec/harness/hashlib; 200-message os.urandom differentials & 0 divergences across 6 confrontations & \code{preuves\_ \allowbreak{}v110/\allowbreak{}independance.csv} \\
\addlinespace[1.5pt]
Comparative benchmark sha256(1,024 B) — AUDITED v0.13.0 (D191) & tls13-01..04 blocks (hand notation, complete operation init+absorb+final) & 75\,983 instructions/op (69.84 instr/round); 765 B code + 256 B table; 27.6 ops/s (emulated) & \code{preuves\_ \allowbreak{}v130/\allowbreak{}audit\_ \allowbreak{}benchmark.csv} \\
\addlinespace[1.5pt]
Comparative benchmark sha256(1,024 B) — AUDITED v0.13.0 (D191) & gcc -O2 reference (same emulator, same 1,024-byte input) & 69\,132 instructions/op (63.54 instr/round); 1\,226 B code + 256 B table; 30.3 ops/s (emulated) & \code{preuves\_ \allowbreak{}v130/\allowbreak{}audit\_ \allowbreak{}benchmark.csv} \\
\addlinespace[1.5pt]
Comparative benchmark sha256(1,024 B) — AUDITED v0.13.0 (D191) & gcc -Os reference (same emulator, same input) & 77\,505 instructions/op (71.24 instr/round); 292 B code + 256 B table; 26.4 ops/s (emulated) & \code{preuves\_ \allowbreak{}v130/\allowbreak{}audit\_ \allowbreak{}benchmark.csv} \\
\addlinespace[1.5pt]
Formal model of CF3 (v0.12.0, D187 — explicit-state exploration) & P1 secret under threshold; P2 roll-back refusal; P3 no destruction on wrong phrase; P4 bounded blast radius & all 4 properties PROVED IN THE MODEL (0 counterexamples); 154 states, 1\,152 transitions, depth 6; adversary at most 2 of 3 elements with epoch replay & \code{preuves\_ \allowbreak{}v120/\allowbreak{}modele\_ \allowbreak{}formel.csv} \\
\addlinespace[1.5pt]
Seal statistics — CORRECTED v0.13.0 (D190) & rule: seal iff $n$ consecutive identical rebuilds, per-rebuild divergence $q$ & P(seal$|q{=}50\%)$: 3.1\% ($n{=}5$), 0.098\% ($n{=}10$); P(seal$|q{=}1\%$): 95.1\%, 90.4\%; one-sided quantiles 1.86$\sigma$/3.10$\sigma$ & \code{preuves\_ \allowbreak{}v130/\allowbreak{}stats\_ \allowbreak{}v130.csv} \\
\addlinespace[1.5pt]
Compiler-reference arm (v0.14.0, D193/D195) & loupe-01 — hand notation vs gcc -O2 C, same contracts, same shims, same vectors & instructions/op: 1025 vs 32 (notation vs compiled); code 448 B vs 494 B; mutation K3 5/16 vs 8/16 — behavior: ZERO divergence on all vectors & \code{preuves\_ \allowbreak{}v140/\allowbreak{}bras\_ \allowbreak{}c.csv} \\
\addlinespace[1.5pt]
Compiler-reference arm (v0.14.0, D193/D195) & loupe-02 — hand notation vs gcc -O2 C, same contracts, same shims, same vectors & instructions/op: 48 vs 47 (notation vs compiled); code 320 B vs 306 B; mutation K3 5/16 vs 9/16 — behavior: ZERO divergence on all vectors & \code{preuves\_ \allowbreak{}v140/\allowbreak{}bras\_ \allowbreak{}c.csv} \\
\addlinespace[1.5pt]
Compiler-reference arm (v0.14.0, D193/D195) & loupe-05 — hand notation vs gcc -O2 C, same contracts, same shims, same vectors & instructions/op: 9089 vs 9146 (notation vs compiled); code 384 B vs 1707 B; mutation K3 6/16 vs 14/16 — behavior: ZERO divergence on all vectors & \code{preuves\_ \allowbreak{}v140/\allowbreak{}bras\_ \allowbreak{}c.csv} \\
\addlinespace[1.5pt]
Compiler-reference arm (v0.14.0, D193/D195) & loupe-08 — hand notation vs gcc -O2 C, same contracts, same shims, same vectors & instructions/op: 641 vs 477 (notation vs compiled); code 320 B vs 351 B; mutation K3 13/16 vs 12/16 — behavior: ZERO divergence on all vectors & \code{preuves\_ \allowbreak{}v140/\allowbreak{}bras\_ \allowbreak{}c.csv} \\
\addlinespace[1.5pt]
Compiler-reference arm (v0.14.0, D193/D195) & loupe-13 — hand notation vs gcc -O2 C, same contracts, same shims, same vectors & instructions/op: 193 vs 69 (notation vs compiled); code 192 B vs 105 B; mutation K3 9/16 vs 13/16 — behavior: ZERO divergence on all vectors & \code{preuves\_ \allowbreak{}v140/\allowbreak{}bras\_ \allowbreak{}c.csv} \\
\addlinespace[1.5pt]
Compiler-reference arm (v0.14.0, D193/D195) & loupe-16 — hand notation vs gcc -O2 C, same contracts, same shims, same vectors & instructions/op: 39 vs 16 (notation vs compiled); code 704 B vs 413 B; mutation K3 9/16 vs 14/16 — behavior: ZERO divergence on all vectors & \code{preuves\_ \allowbreak{}v140/\allowbreak{}bras\_ \allowbreak{}c.csv} \\
\addlinespace[1.5pt]
RQ2 re-measured (v0.14.0, the Loupe-era card) & localiser — end-to-end context tokens (cl100k\_base, n=10), reasoning estimated ($\times$1.08, named) & median 580 [q1 496 .. q3 706] — M1 historical: 51 & \code{preuves\_ \allowbreak{}v140/\allowbreak{}remesure\_ \allowbreak{}rq2.csv} \\
\addlinespace[1.5pt]
RQ2 re-measured (v0.14.0, the Loupe-era card) & éditer — end-to-end context tokens (cl100k\_base, n=10), reasoning estimated ($\times$1.08, named) & median 2,143 [q1 1,085 .. q3 3,463] — M1 historical: 719 & \code{preuves\_ \allowbreak{}v140/\allowbreak{}remesure\_ \allowbreak{}rq2.csv} \\
\addlinespace[1.5pt]
RQ2 re-measured (v0.14.0, the Loupe-era card) & ajouter — end-to-end context tokens (cl100k\_base, n=10), reasoning estimated ($\times$1.08, named) & median 3,510 [q1 3,128 .. q3 4,091] — M1 historical: 1\,088 & \code{preuves\_ \allowbreak{}v140/\allowbreak{}remesure\_ \allowbreak{}rq2.csv} \\
\addlinespace[1.5pt]
RQ2 naive baseline (v0.14.0) & full disassembly of loupe.exe (what the card spares) & 300\,304 tokens (37\,538 disassembly lines) — two orders of magnitude above any card operation & \code{preuves\_ \allowbreak{}v140/\allowbreak{}remesure\_ \allowbreak{}rq2.csv} \\
\addlinespace[1.5pt]
Oracle independence, second agent (v0.14.0) & MD5 T[64] + SHA-256 K[64] regenerated by an isolated agent instance from the mathematical definitions (same-family infrastructure — the gap named) & 0 divergences across 128 constants (agent == plan == definitions); integer-exact cube-root route confirms & \code{preuves\_ \allowbreak{}v140/\allowbreak{}independance\_ \allowbreak{}agent.csv} \\
\addlinespace[1.5pt]
Mutation survivors re-run, composition killers (v0.14.0) & loupe domain — the v0.11.0 mutants re-run with the composition harness (winmain scenario / handshake+record / autonomous sys harness) & 232/288 killed (80.6\%) — 56 survivors named (v0.11.0: 83.4\% overall; now 93.3\%) & \code{preuves\_ \allowbreak{}v140/\allowbreak{}rejeu\_ \allowbreak{}survivants.csv} \\
\addlinespace[1.5pt]
Mutation survivors re-run, composition killers (v0.14.0) & tls13 domain — the v0.11.0 mutants re-run with the composition harness (winmain scenario / handshake+record / autonomous sys harness) & 448/448 killed (100.0\%) — 0 survivors named (v0.11.0: 83.4\% overall; now 93.3\%) & \code{preuves\_ \allowbreak{}v140/\allowbreak{}rejeu\_ \allowbreak{}survivants.csv} \\
\addlinespace[1.5pt]
Mutation survivors re-run, composition killers (v0.14.0) & sys domain — the v0.11.0 mutants re-run with the composition harness (winmain scenario / handshake+record / autonomous sys harness) & 96/96 killed (100.0\%) — 0 survivors named (v0.11.0: 83.4\% overall; now 93.3\%) & \code{preuves\_ \allowbreak{}v140/\allowbreak{}rejeu\_ \allowbreak{}survivants.csv} \\
\addlinespace[1.5pt]
Formal model of CF3, P1 (v0.15.0) & secret below threshold (adversary with $\le 2$ of the three elements never learns the secret) & PROVED UNBOUNDED — the inductive invariant, preservation checked per rule over the finite state shapes; epoch-uniform by translation & \code{preuves\_ \allowbreak{}v150/\allowbreak{}modele\_ \allowbreak{}invariant.csv} \\
\addlinespace[1.5pt]
Formal model of CF3, P2 (v0.15.0) & rollback refusal (the epoch never decreases; replaying a past container-plus-trap pair is refused by name) & PROVED UNBOUNDED — the inductive invariant, preservation checked per rule over the finite state shapes; epoch-uniform by translation & \code{preuves\_ \allowbreak{}v150/\allowbreak{}modele\_ \allowbreak{}invariant.csv} \\
\addlinespace[1.5pt]
Formal model of CF3, P3 (v0.15.0) & no destruction on wrong passphrase (wrong phrases lead only to the benign retry branch) & PROVED UNBOUNDED — the inductive invariant, preservation checked per rule over the finite state shapes; epoch-uniform by translation & \code{preuves\_ \allowbreak{}v150/\allowbreak{}modele\_ \allowbreak{}invariant.csv} \\
\addlinespace[1.5pt]
Formal model of CF3, P4 (v0.15.0) & blast radius bounded to one epoch (the trap media and the CF2 fragments are never touched) & PROVED UNBOUNDED — the inductive invariant, preservation checked per rule over the finite state shapes; epoch-uniform by translation & \code{preuves\_ \allowbreak{}v150/\allowbreak{}modele\_ \allowbreak{}invariant.csv} \\
\addlinespace[1.5pt]
Model-implementation conformance (v0.15.0, directed) & every abstract rule replayed against the Python vault, all failure paths included & 34 steps, 0 gap(s) — five epochs lived at both levels & \code{preuves\_ \allowbreak{}v150/\allowbreak{}raffinement.csv} \\
\addlinespace[1.5pt]
Model-implementation conformance (v0.18.0, generative) & 10\,000 random transition sequences (seed 20261112, 3--8 steps, every rule and refusal path covered), model and concrete vault driven in parallel, ddmin shrinking & 660 divergent traces, ONE class (session falsification: the model's guard is canonical-only, the concrete sentinel destroys on proven falsification mid-session — P3/P4 hold on every divergent trace); 655 shrunk to the 2-step minimum; replay bit-identical (verified twice) & \code{preuves\_ \allowbreak{}v150/\allowbreak{}pbt\_ \allowbreak{}traces.csv} + the v0.18.0 replays \\
\addlinespace[1.5pt]
Coverage lemma + TLA+ alignment (v0.18.0) & the 288 shapes are exactly the abstraction of the well-typed invariant states (bijection verified); the domain constraints inductive (Apalache, unbounded integers, negative control caught); CF3.tla aligned rule by rule with the operational semantics (512 states per rule) & coverage VERIFIED; 13/13 rules aligned; the model-checker battery re-run 8/8 green, verdicts.json identical to the consigned & \code{tla\_ \allowbreak{}v150/} + the v0.18.0 checks \\
\addlinespace[1.5pt]
Wall-clock constant time (hardware dudect kit, 13 Nov 2026, signed witnesses) & four targets from the sealed proxy piece, RDTSC/QPC, core 0 & N=200/class: $|t|$ 12.39, 5.19, 0.94, 3.11 (two above the $>5$ threshold, consigned as measured); N=400/class: 0.84, 1.55, $-0.30$, 2.02 (all under; the dependencies dissipated) — no threshold divergence on that sample, that machine, one run each; never ``constant time proven'' & the signed kit-2 witnesses (Entries 261/264/266) \\
\addlinespace[1.5pt]
RQ2 re-measured (v0.17.0/v0.18.0, the instrumented round-trip) & localiser / éditer / ajouter — the native tokenizer on every logged call, the pre-registered plan ede70035, cl100k in parallel (0.98--1.02) & medians 1\,277 [846--1\,568] / 3\,374 [2\,337--5\,444] / 3\,674 [3\,428--4\,258], $n{=}18/18/15$; the perimeter decomposition confirmed by independent re-tokenization ($n{=}5$ per operation, zero divergence beyond one token) & \code{p4-mesures-brutes.jsonl} + the v0.18.0 decomposition \\
\addlinespace[1.5pt]
RQ1 exact interval (v0.13.0) & 11 of 12 terrain runs surfaced at least one escaped defect & Clopper-Pearson 95\% CI [0.615 — 0.998] (descriptive: sequential iterations of one developer-agent pair) & \code{preuves\_ \allowbreak{}v130/\allowbreak{}stats\_ \allowbreak{}v130.csv} \\
\addlinespace[1.5pt]
\bottomrule
\end{longtable}\endgroup


% ------------------------------------------------------------------
\subsection{The D204 two-arm field campaign}\label{sec:campagne}

The closing validation programme is a pre-registered comparison of the
two ways of working on the same contracts (decision D204, its plan
hash-consigned before any measurement: \code{a78caf8c\ldots}). Each
campaign unit is delivered twice, over one shared layer-I contract:
arm~T --- the full TRIADE discipline (R1--R11: the notation written by
hand, the round-trip, the card, the oracles regenerated from I, the
crystallisation, the guard, the field kit) --- and arm~C --- the
standard compiled flow (the same contract read, the C source written,
gcc~14.2 -O2 cross-compiled for Windows, the tests as the standard flow
practices them, the same field-kit mechanism). The same agent, the
same model family (\code{glm-4-plus}, the identifier reported by every
logged call), the same developer, the same real target machine: the
arms differ only in the discipline. The plan fixes four metrics:
M1 (escaped defects per delivery, under the D202 rule and its named
exclusions), M2 (real tokens, read from the usage logs of every
instrumented call --- the $\times$1.08 estimation forbidden where a log
exists; a missing call is a named gap, never a reconstructed number),
M3 (the developer's human time at the gesture, never recalled from
memory), and M4 (verification compute time, monotone chronometer,
fixed command lists per arm). Directional, falsifiable hypotheses were
registered before the first unit (H1 final health at both arms; H2
escaped defects within a factor of two; H3 the notation arm costlier in
tokens; H4-D204 longer human ceremonies for T; H5 verification time
within one order of magnitude), with stop rules: at most three
deliveries per block, first green closes the unit, an app at most
twelve deliveries, the campaign stops at a red guard or a compromised
piece.

\paragraph{The asymmetric delivery formats, and why.} The two arms do
not deliver the same \emph{binary format}, and the asymmetry is
load-bearing for reading the escapes. Arm~T's piece is an EXE image
composed at a fixed slot of the sealed \code{loupe.exe} layout (the
notation is assembled into the image the card already indexes; the
terrain kit maps it by sections and the scenario calls it at its
address). Arm~C's piece is a DLL (\code{gcc -O2 -shared}, one export
per block declared in a \code{.def} file), because the standard flow's
unit of delivery is a library the host loads: on the terrain the
Windows loader binds the DLL's import table (every leaf must resolve),
while the laboratory's POSIX proof maps the image by sections and
services the imports through shims behind a SysV$\to$msx64 trampoline.
The loading paths differ exactly where the campaign's escapes live:
the unbound-import-leaf class exists only on the DLL path (an EXE
without that import has no leaf to leave unbound), and it struck the
compiled arm three times --- twice at the same slot (the memcpy leaf,
IAT @+0x151a0: the scenario and then the host) and once at the memcmp
leaf (loupe-13, IAT @+0x91a0). Table~\ref{tab:campagne}
reports the complete campaign --- fourteen units (six controlled blocks
and the application, each at both arms), twenty-two witnessed
deliveries, every verdict, every named escape with its attribution
(kit, judge, or piece) --- generated from the ledgers and reconciled
programmatically: 22 deliveries (15 green, 7 red), 6 counted escapes
under the D202 delivery rule, 2 named-and-excluded escapes (named at
the laboratory before any field run), the seventh red delivery
identified as the pre-named judge artifact (a laboratory-judge fault
consumes no delivery), and zero escapes attributable to either arm's
native code. The reconciliation script recomputes every total from the
raw registers and refuses any divergence.

% Généré par reconciliation_campagne_v180.py — ne pas éditer à la main.
\begingroup\scriptsize
\begin{longtable}{@{}l c c l p{4.1cm} l@{}}
\caption{The 22 witnessed deliveries of the D204 campaign, every row generated from the ledgers (livraisons.csv, defauts.csv): unit, arm, iteration, verdict, the named escape if any (exclusion rule D202 applied), and its attribution. The seventh red (C-loupe-05-it1) is the pre-named judge artifact (Entry 294): a laboratory-judge fault consumes no delivery.}\label{tab:campagne}\\
\toprule
\textbf{Delivery} & \textbf{Arm} & \textbf{It.} & \textbf{Verdict} & \textbf{Named escape} & \textbf{Attribution} \\
\midrule
\endhead
\bottomrule
\endfoot
\code{T-loupe-01-it1} & T & 1 & VERT & --- & — \\
\code{C-loupe-01-it1} & C & 1 & VERT & --- & — \\
\code{T-loupe-02-it1} & T & 1 & VERT & --- & — \\
\code{C-loupe-02-it1} & C & 1 & VERT & --- & — \\
\code{T-loupe-05-it1} & T & 1 & VERT & --- & — \\
\code{C-loupe-05-it1} & C & 1 & ROUGE & --- & juge (labo) \\
\code{C-loupe-05-it2} & C & 2 & VERT & --- & — \\
\code{T-loupe-08-it1} & T & 1 & ROUGE & harnais-v3 & kit \\
\code{T-loupe-08-it2} & T & 2 & VERT & --- & — \\
\code{C-loupe-08-it1} & C & 1 & VERT & --- & — \\
\code{T-loupe-13-it1} & T & 1 & VERT & --- & — \\
\code{C-loupe-13-it1} & C & 1 & ROUGE & harnais-feuille-memcmp & kit \\
\code{C-loupe-13-it2} & C & 2 & VERT & --- & — \\
\code{T-loupe-16-it1} & T & 1 & VERT & --- & — \\
\code{C-loupe-16-it1} & C & 1 & VERT & --- & — \\
\code{T-app-it1} & T & 1 & ROUGE & lanceur-prose-non-remmee & kit \\
\code{T-app-it2} & T & 2 & ROUGE & harnais-console-locale & kit \\
\code{C-app-it1} & C & 1 & ROUGE & harnais-feuille-memcpy & kit \\
\code{T-app-it3} & T & 3 & VERT & --- & — \\
\code{C-app-it2} & C & 2 & ROUGE & hote-feuille-memcpy & kit \\
\code{T-app-it4} & T & 4 & VERT & --- & — \\
\code{C-app-it3} & C & 3 & VERT & --- & — \\
\end{longtable}\endgroup


\section{Results}\label{sec:results}

\subsection{RQ1: escaped-defect rate}\label{sec:rq1}

Table~\ref{tab:runs} reports the twelve terrain runs of the Loupe case
study. Eleven of the twelve iterations surfaced at least one escaped
defect --- a defect that had passed every laboratory control in force at
the time, including the three Unicorn harnesses (159 controls at three
fixed seeds), the structural oracle (43 then 121 locks) and the seal
guard. The twelfth run (the 0.1.1 artifact, whose card-debt had been paid
at birth) was green end-to-end. The exact Clopper-Pearson 95\% interval
for eleven of twelve is $[0.615, 0.998]$ --- a \emph{descriptive}
binomial summary only: the twelve iterations are sequential deliveries
of one developer--agent pair, not independent trials, so the interval
must not be read as an inference over a population of developers.
Counting named classes across the
campaign: eleven named defect classes over eleven failing iterations,
each with a named mechanism and an archived evidence piece
(\code{preuves-run-loupe-*} in the evidence vault). The showcase ---
the application simply being \emph{used} on the terrain --- was alive with
zero error-reporter events at runs 10, 11 and 12, after six consecutive
deaths.

The earlier milestones show the same shape at smaller scale: the M2
instrumentation campaign took eight iterations and named eight classes
before the hot-patch loop was proven on a real machine (D43); the M3
proxy needed eight terrain deliveries before the first complete run
(D63), and its killer defects (kernel32 shadow-space violations, a
\code{sub rsp,0x7100} without stack probing) were all of classes the
emulator could not see. We treat this as the study's central empirical
finding: \emph{passing a complete laboratory programme does not transfer
to the terrain}, and the protocol's response is not more laboratory
coverage but a terrain evidence regime in which every escape becomes a
named, falsified class --- each of the eleven Loupe classes is now
guarded by a control or an audit born from its own mechanism (e.g.\ the
msx64 ABI audit born at build from the SSE class of run 4).

The closing campaign sharpens the same finding at the level of the
piece: across 22 witnessed deliveries of the D204 campaign --- six
controlled blocks and the realistic application, at both arms --- the
sealed notation itself escaped zero times; every counted escape
(6 rows in the register, plus 2 named at the laboratory before any
field run and excluded by rule) was kit-side: the rig around the
artifact --- launcher prose, console codecs, unbound import leaves,
one harness constructor --- never a byte of the agent-written native
code. The seventh red delivery traces to the pre-named judge artifact
of the V4 session class (Entry 294, disclosed at both arms); the
unbound-leaf class carries two manifestations at the same slot (the
scenario and then the host) --- the mechanism, named on the artifact's
bytes at Entr\'ees 304/311/312: the unbound IAT slot for memcpy at
image offset +0x151a0 retains its on-disk content, the bare RVA of the
hint/name-table entry (0x152da), and the indirect call through the
slot jumps to that bare RVA --- the crash address is the retained
hint/name RVA (the memcmp leaf of loupe-13 died the same death at
+0x91a0/0x92da). Under the D202 delivery metric --- what the developer
actually received --- the campaign's rate is 6/22; under the piece
metric, zero. The controlled part closed every block at (or after) its
first green; the realistic part consumed seven deliveries to reach
both arms green and signed.

\paragraph{The pre-registered hypotheses, explicit verdicts.} The five
directional hypotheses registered before the first unit, each with its
measured values (read from the ledgers, never recalled): \textbf{H1}
(final health at both arms) \emph{holds} --- both arms closed green
and signed (T-app-it4 and C-app-it3); \textbf{H2} (escaped defects
within a factor of two between arms) \emph{holds} --- 3 counted
escapes at T against 3 at C ($1.0\times$); \textbf{H3} (the notation
arm costlier in tokens) \emph{falsified} --- the compiled arm consumed
$\times$1.27 the tokens of the notation arm on the one instrumented
unit (67{,}475 against 53{,}158), the direction opposite to the
registration; \textbf{H4-D204} (longer human ceremonies for T)
\emph{holds} --- 104 minutes against 17 at the gesture (the timed
ceremony log; $\times$6.1); \textbf{H5} (verification time within one
order of magnitude) \emph{falsified} --- 52.5\,s against 2.2\,s
($\times$23.9). The label disambiguation stands: H4-D193 (the
v0.14.0 compiler comparison, ``testable to the same degree'') is a
different hypothesis from H4-D204.

\subsection{RQ2: token cost per validated change}\label{sec:rq2}

Table~\ref{tab:budgets} shows the M1 instrumentation (51/719/1{,}088 ---
historical, never replaced). The current measurement is the v0.17.0
re-measurement (Table~\ref{tab:preuves}): the three operations measured
as complete instrumented round-trips with the model's native tokenizer
(the usage log of every call, the $\times$1.08 estimation abandoned),
$n{=}18/18/15$ real samples (the card's own blocks and contracts, seed
consigned in the pre-registered plan \code{ede70035\ldots}),
cl100k\_base in parallel (ratio $0.98$--$1.02$ across operations):
locating a block costs a \emph{median} of 1{,}277 tokens [quartiles
846--1{,}568]; editing one 3{,}374 [2{,}337--5{,}444]; adding one 3{,}674
[3{,}428--4{,}258] --- full distributions and wall latencies in the data
archive.

\paragraph{The perimeter note (mandatory before any comparison).}
The earlier re-measurement (v0.14.0, D195) measured the \emph{content
alone} with cl100k\_base, reasoning estimated ($\times$1.08, named):
medians of 580 [496--706] to locate, 2{,}143 [1{,}085--3{,}463] to edit,
3{,}510 [3{,}128--4{,}091] to add ($n{=}10$) --- historical values,
explicitly labelled as such. The v0.17.0 perimeter is the full
round-trip. The decomposition, generated from the raw logs and
confirmed by an independent re-tokenization of the prompts and responses
($n{=}5$ per operation, zero divergence beyond one token): the locating
operation decomposes into a median of 596 tokens of material read (the
card header --- the part v0.14.0 counted), 41 tokens of request prose,
and 604 tokens of \emph{written response} (the agent re-emits the full
header), all under the native tokenizer; editing decomposes into
1{,}617/45/1{,}635 and adding into 1{,}416/55/2{,}378. (A
methodological note the present edition adds: the component medians are
computed independently per component and do \emph{not} sum to the
operation medians --- $596{+}41{+}604 = 1{,}241$ against 1{,}277 for
locating, $1{,}617{+}45{+}1{,}635 = 3{,}297$ against 3{,}374 for
editing, $1{,}416{+}55{+}2{,}378 = 3{,}849$ against 3{,}674 for adding;
medians are not additive --- the components rank blocks differently ---
and the add-operation's component sum even exceeds its total median.
The totals are the measured round-trip medians, never sums of component
medians.) The apparent
$2.2\times$ growth in locating (580 $\to$ 1{,}277) is therefore a
perimeter effect --- the tokenizer ratio is 0.98--1.02 and cannot
produce it --- and the two values are never compared without this note.
The M1 values (51/719/1{,}088) remain single-shot history; the
Loupe-era content values cost 2.98--11.37$\times$ their M1 counterparts
(the notations longer, the contracts denser --- the cost grew with the
rigour; every ratio recalculated from the raw CSVs by the numeric
linter of the present edition, which refuses any divergence).

\paragraph{The baselines.}
What the card-fragment indexing buys is measured against three
baselines, all on the same eighteen blocks. The naive one: reading the
full disassembly of \code{loupe.exe} costs 300{,}304 tokens
(37{,}538 lines) --- 235$\times$ the cheapest card operation (locating,
1{,}277) and 82$\times$ the most expensive (adding, 3{,}674) at the
current medians. The two realistic ones, measured to close the gap
between the card and the naive floor: (a)~the function-level
disassembly by symbol table --- the unit a standard tool would feed the
agent --- costs a median of 2{,}074 tokens [quartiles 1{,}192--3{,}202;
min 317, max 6{,}552]; (b)~a hierarchical retrieval over the disassembly
(the block header plus a fifteen-line window) costs a median of 792
tokens [628--982; min 525, max 1{,}217]. Read at the comparable perimeter
(content, cl100k\_base): the card's locating \emph{material} (a median
of 596) sits below the retrieval baseline (792) --- the indexed fragment
is smaller than header-plus-window; the full round-trips (1{,}277 native)
sit above it, the written response being the difference. The whole
generation-1 lineage totalled 13{,}873 tokens (D34); full-card audits
remain exceptional, logged operations.

\paragraph{The campaign's instrumented unit, decomposed.}
The campaign's one fully instrumented head-to-head unit (loupe-01, both
arms, every call logged) reads 53{,}158 tokens for the notation arm
against 67{,}475 for the compiled arm ($\times$1.27, the direction
opposite to the pre-registered H3, $n{=}1$). The decomposition by
category, read from the usage logs: at arm T --- specification 5{,}313,
oracle 3{,}752, notation 11{,}480, corrective iterations 32{,}613
(61.4\% of the total, six calls); at arm C --- specification 5{,}605,
oracle 4{,}423, C source 14{,}660, test suite 3{,}563, corrective
iterations 39{,}224 (58.1\%, seven calls). Two readings: the corrective
spiral dominates both arms (the hot path is the workflow, not the
notation), and the compiled arm's \emph{source} costs more tokens than
the hand-written notation (14{,}660 against 11{,}480) --- the first
counter-intuition behind H3's fall. The reasoning tokens are not
separated by the API log (counted within the completion): a named gap,
never estimated back.

\paragraph{The verification-time metric (M4) and its optimization.}
The campaign's verification-time metric reads M4 $=$ 52.5\,s (arm~T)
against 2.2\,s (arm~C) on the instrumented units --- $23.9\times$,
falsifying the pre-registered H5 (``within one order of magnitude''),
$n{=}1$, stated as such. The present edition optimizes the arm-T
fixed command list \emph{without removing any control} (rule R11: no
optimization without measurement): the twelve crystallisation rebuilds
run concurrently (each in a hermetic copy of the unit folder, all
outputs compared byte for byte --- determinism is order-independent);
the oracle differentials run as parallel sessions, each file in its own
persistent Unicorn emulation (the profile: the session \emph{opening}
--- the emulated md5+sha256 over the whole file --- dominates at
3.2\,s per 65\,KiB, the reads cost 0.15\,ms each). Measured
before/after on the same two-core laboratory machine (monotonic clock,
the campaign's fixed command list): the before-measure reads 50.9\,s
--- 3.0\% below the campaign's consigned 52.5\,s, the run-to-run
spread of two single measurements ($n{=}1$ each, both consigned as
measured, never averaged into an invented baseline) --- and the
after-measure is two consecutive runs ($n{=}2$), 39.3\,s then
39.2\,s, both green, sealed fingerprints unchanged, and the oracle
consoles \emph{identical line for line} to the originals: a $23\%$
reduction, with
the compiled arm re-measured on the same machine at 2.2\,s --- the
ratio moves from $23.1\times$ to $17.9\times$ and H5 \emph{remains
falsified}; the remaining critical path is the emulated cryptography of
the contract sessions themselves ($\approx$10\,s for the 196\,KiB
witness), irreducible without changing the controls. The coverage
compromise refused: sharing one emulation between the build, the
round-trip and the oracle was considered and rejected --- the judge
never shares state with the judged.

\subsection{RQ3: defect classes invisible to emulation}\label{sec:rq3}

The eleven named classes of the Loupe campaign group into five families;
each was invisible to the emulation harness for a structural reason, not
a coverage gap --- and each corresponds to a guarantee a compiler
normally provides \emph{silently} (a loadable, relocatable image;
ABI-conformant call sequences; conventional error codes; runtime
assumptions baked into the toolchain), which hand notation must
re-derive explicitly, one named lesson at a time:

\begin{enumerate}
\item \emph{Loader and image semantics} (runs 1--3): an ASLR image
without a relocation table (refusal \code{0xc000\allowbreak 007b}); an image outside
the proven envelope of eight prior artifacts (page hole, massive BSS);
an uninitialized \code{.data} session frame (the initializer existed only
in the harness's ideal memory). The harness parses images; it never loads
them.
\item \emph{ABI conventions with hardware consequences} (run 4): msx64
stack arguments 5+ written at the home space instead of
\code{[rsp+0x20]+}, and stack misalignment at call sites
($\mathrm{rsp} \equiv 8 \bmod 16$) --- harmless under emulation, fatal
with a real SSE-using \code{gdi32}. This class gave birth to a static
msx64 ABI audit at build time, which then caught two further alignment
defects no one had seen.
\item \emph{Operating-system loader behaviour} (runs 5, 9, 11): Windows
refusing a write-open on the file of a live mapped section
(\code{gle 32}); import tables walked by the loader at \code{DWORD}
stride while the build wrote \code{QWORD} entries --- the first cause of
all six showcase deaths, found on the artifact's bytes; and the
structural false negative that a library loader never resolves an EXE's
imports (the ``liaison'' probe could never succeed as designed --- the
showcase itself is the proof).
\item \emph{Error-code conventions} (runs 6--7): an \code{hInstance}
read from the PEB instead of \code{GetModule\allowbreak HandleW\allowbreak (NULL)}; negative
return codes read as unsigned 32-bit values by the kit while the piece
delivered them sign-extended.
\item \emph{Harness- and kit-embedded assumptions} (runs 8, 10 and the
pre-run fix of run 11): a probe killed by its own unguarded Python; a
64-bit \code{HMODULE} truncated by \code{ctypes} to a C \code{int}
because \code{argtypes} was not set; a bare \code{ret} neutralizing an
entry point whose return value the loader reads in \code{DllMain}
semantics. These are defects of the \emph{test rig}, named with the same
discipline as defects of the artifact.
\end{enumerate}

The pattern across families: the emulator verifies \emph{logic} well and
\emph{environment} not at all. The protocol's answer is asymmetric ---
add the static audits that can see the environment-sensitive classes
(ABI, image format), and route the rest through terrain evidence with
named-defect accounting, rather than pretending the harness can be
completed.

\subsection{RQ4: hot-patch safety}\label{sec:rq4}

The M2 campaign demonstrated the full loop on a real machine: replacing
the echo block \emph{while executing}, verifying by comparison,
rolling back, re-verifying --- with zero post-patch regression across the
eight instrumented iterations (D43). On RISC-V (M7), where fixed-width
instructions forbid the x86-64 relative trampoline, the patch path uses
a scratch register plus \code{fence.i}. The M7 terrain evidence for this
path was produced on the developer's real machine (WSL2, a real kernel
and real syscalls) executing the RISC-V artifact under the kit's
embarked \code{qemu-riscv64} process emulator --- the instrumental
reserve named in every M7 verdict (D149): the ISA was emulated, never
disguised as native execution. The patch \emph{mechanism} is proven in
that regime (notation, round-trip, the in-traffic switch with byte-exact
transcription, falsification, zero regression; 29/29 recoup on crossed
pieces); the hardware instruction-cache semantics of \code{fence.i} on
physical RISC-V silicon remain \emph{unexecuted} and are explicitly not
claimed. Load-time verification (R10)
checks per-block checksums and signed version monotonicity; the
anti-rollback barrier was exercised at each sealing (reconstruction of a
sealed version is refused, observed at the sealing gesture itself), and
version epochs open strictly monotonically --- the 0.1.1 artifact
$(0,1,1,-1,1)$ opened the new epoch; a pre-release of the sealed 0.1.0
remains refused. Sealed artifacts never changed after their sealing run
(Table~\ref{tab:empreintes}); the 0.1.1 release refreshed a stale
checksum table \emph{at birth} --- a debt named at the 0.1.0 sealing
(D178): the table lives \emph{inside the artifact image} (the R10
load-time directory, seventeen per-block sums plus the version qword),
and the piece-side retouches consigned at Entr\'ees 210/214 had changed
block bytes without refreshing the sums; the recomposed 0.1.1 image
rewrites exactly those seventeen sum fields and the version qword
($0x00010008 \to 0x00010101$), \emph{zero bytes elsewhere} --- which
is why the 0.1.1 image hash differs from 0.1.0
(Table~\ref{tab:empreintes}) while every code block stayed byte-identical
(R7), the sealed 0.1.0 artifact remaining frozen at its own fingerprint
--- an explicit debt-repayment pattern the protocol names and tracks.

% ------------------------------------------------------------------
\subsection{Strengthened evidence (the v0.11.0 programme)}\label{sec:renforce}

The second edition added six laboratory instruments to the evidence base
(Table~\ref{tab:preuves}); all run with fixed, consigned seeds and are
reproducible from the scripts shipped with this edition.

\paragraph{Mutation testing --- re-run with composition killers (v0.14.0).} The 832
single-byte mutants of v0.11.0 were re-run (same seed, same module,
same generation order) with a new semantic killer: a \emph{composition
harness per mutant} --- the full application loop for the Loupe domain
(open, cascade, paint, type, scroll, save, quit; the saved disk and
screen bytes confronted to the sane baseline), the ClientHello composition
plus record protection for the tls13 orchestration blocks, and the
autonomous sys harness for the sys domain (D195). Result: 776/832
killed (93.3\%, up from 83.4\%) --- tls13 and sys now at 100\%, the
56 remaining Loupe survivors \emph{named} (mutations in paths the
scenario does not exercise: error branches, unreached offsets,
neutral bytes). The original campaign: 832 single-byte mutants of the sealed blocks
(instructions, operands, offsets --- one byte mutated in a laboratory copy,
the sealed artifact untouched), killed by the existing instruments in
order of cost: the round-trip check, the structural oracle, and --- for the fourteen
blocks of the cryptographic families (the thirteen tls13 crypto blocks
and loupe-06's md5; 224 mutants, Table~\ref{tab:preuves}) --- the
semantic differential against official vectors. Result: 694/832 killed (83.4\%);
the cryptographic families reach 100\%. The 138 survivors concentrate in
the orchestration and UI blocks, where the semantic killer is not wired
at the isolated-mutant level (their semantics are verified by the full
harnesses in composition, not per mutant): this gap is \emph{named}, not
hidden, and is the primary target of the next edition.

\paragraph{Coverage-guided fuzzing.} The two untrusted-input blocks were
attacked under instrumentation: \code{loupe-\allowbreak 01} (the file path, through
the harness's virtual OS with canaries around the path buffer) reached
60.5\,\% block-instruction coverage over 400 iterations with zero crashes
and zero canary violations; the \code{http-10} parser (on the sealed 0.4.0
artifact) reached 358 distinct instructions with zero emulation crashes
and zero named faults, growing a 17-entry corpus. The limits are stated:
coverage is per-block under emulation, and crashes are emulation faults,
not terrain deaths.

\paragraph{Bounded static memory safety.} An abstract interpretation of
all 52 blocks (the extension of the msx64 ABI audit born at run 4):
2{,}098 memory accesses classified --- 1{,}181 stack accesses verified
against the frame bound (zero out-of-frame: the stack-smash class is
absent), 263 accesses mediated by contract-bounded argument pointers, 191
constant \code{[rip]} accesses, 463 unresolved accesses \emph{named}
(volatile-recomputed bases, composed indices). The analysis is linear
(no path enumeration, loop bodies verified once, no alias analysis): a
bounded verification, not a proof.

\paragraph{Constant-time (emulation --- reformulated as exact trace equality).}
The earlier editions reported the Welch statistics of a t-test on
executed-instruction counts between fixed-secret and random-secret
classes (N=60 each) as uniformly zero. That formulation was an artefact
twice over, both
facets found by the present edition's replay. First: on bit-identical
instruction traces the variance is null and Welch's $t$ is \emph{undefined}
(0/0), not zero --- a t-test is the wrong instrument for an equality that
is exact. Second: on two of the six targets (hmac, aes-gcm) the v0.11.0
instrument had never activated its instruction counter (the counts were
all zeros; the hkdf target had run with an accidentally truthy flag) ---
the vacuous measurements are named, and the replay measures all six for
real. The correct statement is \emph{exact equality of the executed-
instruction traces}, and the replay verifies it as such: each of the six
targets (sha-256, hmac, hkdf, x25519, aes-gcm, md5) executes a single
instruction count across all 120 measurements --- sha-256 4{,}469, hmac
18{,}371, hkdf 18{,}407, x25519 exactly 1{,}238{,}959 for every key,
aes-gcm 11{,}695, md5 132{,}509 (grain consigned, the v0.11.0 protocols
replayed verbatim). N=60 per class is an exact \emph{structural}
control, never a statistical power test. The metric remains
emulation-structural; the real-hardware wall-clock timing is a separate,
now measured, claim (the dudect kit runs below).

\paragraph{Comparative benchmarks --- audited and corrected (v0.13.0).} The
v0.11.0 edition reported a ``15.3$\times$ instruction advantage'' for
the hand-written SHA-256 blocks over a compiled reference. The audit
required by the present edition found that figure to be a \emph{measurement
artifact}: the instrument read the emulation counter after the hash
pipeline returned, but the counter resets at every block invocation ---
the published 4{,}524 instructions were the \emph{finalisation stage
alone}, compared against the reference's \emph{complete} operation (the
implausible 4.2 instructions per SHA-256 round was the tell; a scalar
round needs an order of magnitude more). Re-measured stage by stage,
digest verified against \code{hashlib} before counting and the 17
compression invocations confirmed mechanically, the complete operation
costs 75{,}983 instructions for the blocks against 69{,}132 for
gcc~14.2 -O2 and 77{,}505 for -Os, in the same emulator on the same
input --- under the explicit convention that every percentage here is
\emph{the cost of the notation relative to the reference, positive =
the notation more expensive}: the hand notation sits \emph{between}
the two compiler levels at $+9.9\%$ executed instructions against -O2
and $-2.0\%$ against -Os (the v0.11.0 text had both signs inverted;
every sign and value is recalculated from the raw CSV by the numeric
linter of the present edition, which now checks the signs too), is
1.45$\times$ more compact than the -O2 function (1{,}021 vs 1{,}482
bytes including the shared K table) but less compact than -Os (548
bytes), and neither camp uses any ISA extension beyond the x86-64
baseline --- no AVX, no SHA-NI (the -O2 reference does use baseline SSE2
moves for buffer setup, which the pure-GPR notation does not; the gap
is reported unadjusted and strict GPR-to-GPR parity is not claimed).
The 15.3$\times$ claim is \emph{withdrawn}; the corrected statement is
that the WYSIWYG notation is competitive with a modern compiler on its
own turf --- the price of direct auditability is nearly nil, and the
spectacular premium does not exist.

\paragraph{Oracle independence --- second agent (v0.14.0).} The oracle was
re-generated by a second, \emph{isolated} agent instance (fresh context,
no shared conversation) deriving the MD5 $T[64]$ and SHA-256 $K[64]$
tables from the mathematical definitions alone: \emph{zero divergences}
across the 128 constants (agent == plan == definitions), confirmed by an
integer-exact cube-root route executed independently. The named gap:
the laboratory's agent infrastructure exposes only one model family ---
true cross-family independence remains open. The v0.11.0 routes: the
MD5/SHA tables re-derived
through decimal arithmetic, the RFC 1321 gold vectors confronted
spec$\leftrightarrow$\allowbreak harness$\leftrightarrow$\allowbreak hashlib, the differentials
re-run with a new seed and the system generator (\code{os.urandom}):
zero divergences across all confrontations. A second \emph{agent} remains
to be tried (the replication protocol is proposed in the next edition).

% ------------------------------------------------------------------
\subsection{The controlled compiler comparison (the v0.14.0 programme)}\label{sec:comparaison}

The pre-registered plan (D193, hash-consigned before any measurement)
compares six sealed Loupe blocks spanning the five families against a C
reference implemented \emph{from the same layer-I contracts}, compiled
with gcc 14.2 -O2/-Os (clang absent, the deviation named), calling the
\emph{same mini-OS shims} through an \code{ms\_abi} pointer table, and
confronted with the \emph{same vectors} (D195). The headline result is
\emph{behavioural}: zero divergence across all contract vectors ---
the hand notation and the compiled reference are observationally
equivalent at the contract level.

\emph{Instructions (complete operations, the audited counting convention).}
On SHA-256 range hashing the notation is at \emph{parity at equal work}:
9{,}089 vs 9{,}146 instructions (the 384-byte block composes the sealed
765-byte inherited implementation through the card; the 1{,}707-byte C
function embeds its own). Windowed file reads are at parity (48 vs 47). On
the simpler pure blocks the compiler's tight loops win: hex-line
composition 641 vs 477, byte search 193 vs 69, keystroke entry 39 vs 16.
The full session-open cascade (window + fingerprint panel) reads
102{,}176 vs 75{,}560 --- \emph{not at equal work}: the compiled camp
computes only the SHA-256 half of the panel (no audited C MD5 exists;
the v0.11.0 reference was rejected as faulty), so the comparison
carries its mark until an audited C MD5 exists. The per-vector audit
of the present edition names a further artefact: the loupe-01 median
(1{,}025 vs 32 instructions) falls on a \emph{refusal} vector (four
of the five suite vectors are refusals), where the notation composes
its controls in-piece and the compiled reference stays short --- the
real opening costs 102{,}136 vs 75{,}525 ($1.35\times$, still not at
equal work), reported per vector in the data archive.

\emph{Sizes.} Mixed, honestly: the notation is more compact on
SHA-256 (384+765 shared inheritance bytes vs 1{,}707) and on the
session-open block (448 vs 494 bytes), less compact on keystroke
entry (704 vs 413) and search (192 vs 105) --- the directions
corrected in the present edition against the raw table, where the
previous edition had inverted the open and keystroke entries.

\emph{Mutation (same protocol, same seed, same draw for both camps).}
Camp notation: 5--13/16 killed per block under the differential
killer at -O2 (3--13 across both optimisation levels; every bound
recalculated from the raw table by the numeric linter of the present
edition), the survivors named (consistent with the orchestration
class). Camp compiled: 16/16 on the OS-calling blocks (the call-
sequence killer catches everything), 12--14/16 under the differential
killer on the pure blocks --- and 16/16 in \emph{total} everywhere,
the remaining mutants falling to the cheaper killers. The reading we
pre-registered as H4-D193 holds at the level that matters: both camps
kill 96/96 mutants in total --- \emph{testable to the same degree},
the killer profiles differing (the notation's kills concentrate in
the differential killer, the compiled camp's spread across
differential and call-sequence).


\subsection{The unbounded model and conformance testing
(the v0.15.0 programme)}\label{sec:v150}

\paragraph{From bounded to unbounded.} The v0.12.0 model checked the
vault's claims by exhaustive exploration up to three epochs --- an
artificial ceiling inherited from the explorer, not from the design (the
implementation's closing step is $e \leftarrow e+1$ without a cap). The
v0.15.0 programme removes the ceiling and moves the proof to an
\emph{inductive invariant} $I$ over the operational semantics:
$I_1$ (a learned secret implies the full three-element threshold),
$I_2$ (no unjust destruction flag), $I_3$ (the trap media and the CF2
fragments are never touched), $I_4$ (epochs stay positive). The proof
obligations are discharged mechanically: initialisation holds at the
initial state; preservation is checked for \emph{each of the 13 rules}
over the finite set of state shapes (288 shapes satisfying $I$, the
epoch instantiated over both guard classes $\{1\}$ and $\{\ge 2\}$); the
epoch-uniformity of the class $\ge 2$ is \emph{structural} --- the
rules only ever compare the epoch to 1 and take its successor, a fact
of the rule text audited mechanically (the epoch is read by
$\{e > 1, e + 1\}$ only) --- and witnessed in replay by re-running the
whole proof under epoch translations $\delta \in \{7, 50, 1000\}$ with
bit-identical verdicts. The present edition adds the \emph{coverage
lemma} that was previously implicit: every concrete reachable state is
represented by at least one of the 288 abstract shapes, so the
preservation check over the shapes covers the entire reachable space.
The lemma is discharged in three obligations: (O1)~the abstraction
$\alpha$ is total on well-typed states (its domains are the typing
constraints); (O2)~those domain constraints are \emph{inductive} ---
verified by Apalache over unbounded integers (initialization and
preservation both NoError, with a negative control: a weakened domain
is caught, Error~12); (O3)~the 288 shapes are \emph{exactly} the image
of the well-typed invariant-satisfying states under $\alpha$ --- a
verified bijection (288~forms, none outside the invariant, all
reachable bounded states covered). The TLA+ model submitted to the
model checkers is aligned with the operational semantics \emph{rule by
rule}: a strict transpiler reads the thirteen actions of \code{CF3.tla}
and confronts guards and effects with the Python semantics over 512
concrete states per rule (both epoch classes, every subset, every
boolean) --- all thirteen aligned, any unrecognized construct a refusal.
The independent model-checker replay of the present edition (TLC the
bounded fidelity: 154 states, 1{,}152 transitions reproduced exactly;
Apalache/SMT the unbounded induction
$\mathrm{Init} \Rightarrow I$ and $I \wedge \mathrm{Next} \Rightarrow I'$,
negative controls caught) returns 8/8 green. Monotonicity (no rule ever decreases the epoch)
is checked unconditionally over 4{,}096 shapes; the two replay rules
return the state unchanged. The four properties are therefore
\emph{proved unbounded}: any epoch, any depth, under hypotheses
A1--A4, with A5 (side channels, implementation flaws, three-element
compromise, quantum) explicitly out of model.

\paragraph{Fidelity and bounded cross-checks.} The ported explorer
reproduces the v0.12.0 counts \emph{exactly} under the historical
ceiling (154 states, 1{,}152 transitions --- the proof of a faithful
port); with the ceiling removed, the extended bounds ($e \le 5$ at depth
8: 228 states / 1{,}902 transitions; $e \le 7$ at depth 12: 372 /
3{,}342) find zero reachable states outside $I$ and zero property
counterexamples --- the bounded evidence now confirms the unbounded
proof instead of substituting for it.

\paragraph{Model--implementation conformance (renamed in the present
edition; formerly ``refinement'') --- now a generative campaign.} The
v0.15.0 evidence was a directed game of 34 replayed steps (every
abstract rule against the real Python implementation, a test vault
under a publicly marked test passphrase; the real vault is never opened
by a test): the honest cycle across \emph{five} lived epochs, fifteen
wrong-passphrase attempts (zero destruction, zero sealing), the proven
container falsification, the trap falsification, the past-epoch replay
in three variants, the absent trap, the foreign vault's trap, and the
below-threshold adversary (32 forged key-halves triples, zero verifier
passes) --- zero gaps, every failure path included. The present edition
replaces the chosen-steps design with a \emph{generative} campaign:
10{,}000 random transition sequences (seed consigned, 3--8 steps per
trace, weighted draws over the actions \emph{offered} by the model,
every rule and every refusal path covered --- accepted and refused
counts per rule are read from the generation logs), each trace driving
the abstract model and the concrete test vault in parallel, every
verdict confronted step by step, every divergence shrunk (ddmin) and
archived. Result: \textbf{660 divergent traces, exactly one class} ---
\code{falsifie\_conteneur} during an open session: the model's guard
offers the rule only in the canonical phase, while the concrete
sentinel runs on the canonical container and \emph{does} destroy on
proven falsification mid-session (local destruction, PV raised, the
trap media and CF2 untouched). The analysis: a state-coverage gap of
the \emph{model} (an under-approximation discovered by the campaign
--- the directed 34-step game never crossed it, the randomization
crosses it 660 times in 10{,}000), \emph{not} a threat to the
properties: the concrete destruction remains proven falsification and
remains local (P3 and P4 hold on every divergent trace); the shrinker
(ddmin) reduces every divergent trace to its minimal failing sequence:
655 of the 660 reach the two-step minimum [\code{ouvrir\_correct},
\code{falsifie\_conteneur}]; the five remaining traces --- of the same
single class --- shrink to minimal sequences longer than two steps,
their per-trace minima consigned in the campaign's raw register (the
655/660 count is read from that register; the analysed object is the
class, not the trace length). The replay at
identical seed reproduces the 10{,}000 verdict digests bit for
bit (verified twice in the present edition, plus the consigned
original). Three concrete gestures sit outside the abstract model
(the out-of-band epoch witness, the foreign vault, the forged halves)
and remain consigned as such, with their concrete verdicts matching
the specification.

\paragraph{The hardware runs: the dudect timing measured; two kits still
waiting.} The real-time dudect kit ran on the developer's machine,
twice, with dated and hand-signed witnesses (13 November 2026; four
targets extracted from the sealed proxy piece, RDTSC/QPC timing, core
affinity 0, the native call validated against the FIPS~180-4 ``abc''
vector). At N=200 per class the Welch statistics read $|t| = 12.39$
(tls13-02) and $5.19$ (tls13-13) --- above the $|t|>5$ threshold,
consigned exactly as measured --- with tls13-16 at $0.94$ and tls13-17
at $3.11$; the N=400 run reads $0.84$, $1.55$, $-0.30$ and $2.02$ ---
all four under the threshold, the N=200 dependencies dissipated at the
larger sample. The honest reading is the kit's own: no threshold
divergence on that sample, that machine, one run at each N --- never
``constant time proven''; cache-timing and power side channels remain
out of scope, and the single-run scope is part of the claim. Two kits
remain \emph{awaiting a hardware run}: \code{fence.i} semantics on
physical RISC-V silicon (the lived M7 terrain runs were TCG emulation
--- the reserve is named), and a third-machine Loupe field run (the
full T0--T7+S0 scenario). For those two, no number exists and none is
fabricated; the claims table stays honest until the runs are on record.

% Table des statuts systématiques — v0.18.1 : rendue en longtable par
% l'audit de la vague 6 (instruction bqh4qlrqj, point 10) : la table
% flottante de la v0.18.0 était TRONQUÉE à la dernière page du PDF
% (les lignes « Compiler-reference comparison » — partielle —,
% « Campaign D204 » et « Independent replication » étaient perdues).
\begingroup\scriptsize
\begin{longtable}{@{}p{3.4cm} p{2.9cm} p{8.4cm}@{}}
\caption{Systematic proof statuses (the v0.15.0 requirement, carried and updated by every edition since): every claim carries its status and its evidence piece.}\label{tab:statuts}\\
\toprule
\textbf{Claim} & \textbf{Status} & \textbf{Evidence} \\
\midrule
\endhead
\bottomrule
\endfoot
P1 secret below threshold & proved unbounded & inductive invariant $I_1$ (secret $\Rightarrow$ threshold 3); preservation by all 13 rules over 288 state shapes; epoch translation $\delta\in\{7,50,1000\}$ uniform \\
P2 rollback refusal & proved unbounded & unconditional epoch monotonicity (4{,}096 shapes, zero decrease); replay rules return the state unchanged \\
P3 no destruction on wrong phrase & proved unbounded & inductive invariant $I_2$: no rule ever writes unjust destruction \\
P4 blast radius one epoch & proved unbounded & inductive invariant $I_3$: no rule ever touches the trap media or the CF2 fragments \\
Coverage of the 288 shapes & proved (v0.18.0) & every reachable concrete state maps to one of the 288 forms: abstraction totality + domain inductivity (Apalache, unbounded integers, negative control caught) + verified bijection with the well-typed invariant states \\
TLA+ $\leftrightarrow$ operational semantics & verified (v0.18.0) & the 13 actions of CF3.tla transpiled strictly and confronted with the Python semantics over 512 states per rule --- all aligned; the model-checker battery re-run 8/8 green, verdicts identical to the consigned \\
Model $\leftrightarrow$ implementation & verified in the laboratory & directed: 34 conformance steps, 0 gaps; generative (v0.18.0): 10{,}000 traces, seed consigned, every rule and refusal path covered --- 660 divergences of ONE class (the model's session-falsification gap; P3/P4 hold on every one), 655/660 shrunk to the 2-step minimum (the five others, of the same class, to longer minima; the raw register consigned), replay bit-identical \\
Bounded cross-checks & proved bounded & the v0.12.0 explorer reproduced exactly (154 states / 1{,}152 transitions); extended bounds 228/1{,}902 ($e\le5$, $d\le8$) and 372/3{,}342 ($e\le7$, $d\le12$), zero states outside $I$ \\
Structural constant time & proved bounded & exact equality of emulated instruction traces across secret classes (single count per target over 120 measurements; $t$ undefined, the ``$|t|{=}0$'' formulation withdrawn) \\
Wall-clock constant time & measured (2 runs, 1 machine) & the dudect kit on the developer's machine, signed witnesses: N=200 --- two targets above the threshold (12.39, 5.19, consigned as measured); N=400 --- all four under (0.84, 1.55, $-0.30$, 2.02); never ``constant time proven'' \\
\code{fence.i} on silicon & awaiting hardware run & the M7 scenario retargeted at a physical RV64 board; the lived M7 runs were TCG emulation, the reserve named --- instruction-cache semantics unproven \\
Third-machine field liveness & awaiting hardware run & the Loupe kit on a second Windows machine, full T0--T7+S0 scenario \\
Compiler-reference comparison & measured (v0.14.0) & zero divergence, SHA-256 parity at equal work; H4 holds \\
Campaign D204 & measured & 22 deliveries reconciled; H1/H2/H4-D204 hold, H3/H5 falsified ($n{=}1$ pairs, stated as such); the M4 optimization measured before/after with no control removed (52.5\,s consigned at the campaign, 50.9\,s the optimization's before-measure, $n{=}2$ after: 39.3/39.2\,s; $23.1\times \to 17.9\times$ on the lab machine, H5 still falsified) \\
Independent replication & awaiting (recruitment / real command) & the T/X grid finalised, the $\kappa$ calculator autotested (synthetic cases), criteria pre-registered unchanged; the extension plan D210 pre-registers $n\ge6$ instrumented units, identical formats, randomized arm order \\
\end{longtable}\endgroup


% ------------------------------------------------------------------
\section{Discussion}\label{sec:discussion}

\paragraph{What the data supports.} A single developer with an LLM agent
delivered, in 17 days, a cross-platform native codebase with a pure
TLS~1.3 stack and a real Windows application, under a discipline that
made every claim checkable: every defect named, every fix diffed, every
seal fingerprinted, every decision append-only. The escaped-defect data
(RQ1/RQ3) is, to our knowledge, the first quantified account of
\emph{where} agent-written native code fails when laboratory verification
is strong but environmental fidelity is absent: not in logic, but in
loader semantics, ABI conventions and rig assumptions. Agent-failure
quantification exists at source level in repository-scale
settings~\cite{yang2024sweagent}, and the verified-compilation line
quantifies the translator's own guarantees~\cite{leroycompcert}; the
claim is scoped to what is new here --- the sealed-notation level, under
signed-intent governance.

\paragraph{What the data does not support.} The study cannot claim that
the protocol is \emph{more} productive or \emph{more} correct than a
compiler-based workflow: one case, one agent family, one developer ---
though the D204 campaign now provides the two-arm comparison on fourteen
units (behaviourally equivalent at the contracts; the cost and
verification-time pairs still $n{=}1$ each). The token budgets were
measured three times --- M1 single-shot, the v0.14.0 content-only
re-measurement, the v0.17.0 instrumented round-trip (the perimeters
documented; never compared without the note) --- and the campaign adds
the only fully instrumented head-to-head unit. The oracle is derived
from the same layer I it checks, so
specification errors are shared, not caught --- mitigated, not removed,
by the second, isolated agent that regenerated the crypto constants
with zero divergences (v0.14.0; cross-family independence remains the
named gap). Memory safety is enforced
only to the depth of the ABI audits and the fuzzing campaigns;
constant-time is verified \emph{structurally} (secret-independent
instruction traces in emulation) --- real-hardware timing remains
unmeasured. These limits are stated, not smoothed over; the companion
programmes of the two subsequent editions addressed several of them
directly (mutation testing, bounded memory safety and structural
constant-time in v0.11.0; the formal model of the vault chain in
v0.12.0), and the benchmark audit of the present edition is itself an
instance of the method applied to its own claims.

\paragraph{On the ``no compiler'' position.} We phrase the protocol's
position precisely: TRIADE eliminates the \emph{optimizing and opaque}
translation chain by having the agent write the assembly notation itself,
with a round-trip check making the notation WYSIWYG
(what-you-write-is-what-executes) at every build. It does not eliminate
the assembler, the operating system, or the platform ABI --- those are
called through their documented conventions, like any native binary.
Similarly, the sovereignty programme's claim about expropriation is
narrow: the \emph{evolution} of stolen code is blocked (signed layer-I
intent, sealed inheritance, offline keys, dead-man's-switch Shamir
escrow); the \emph{execution} of a stolen copy is not, and the
specification says so.

% ------------------------------------------------------------------
\subsection{Formal model of CF3, structured threat model, and
replication protocol (the v0.12.0 programme)}\label{sec:formel}

\paragraph{A bounded formal model of the CF3 vault.} The sealed-vault
chain (CF3: three-element lock, two-act cycle, epoch anti-rollback,
invocation-based sentinel, destruction on proven falsification) was
modelled as an explicit-state transition system and model-checked
exhaustively (154 states, 1{,}152 transitions, depth 6, three epochs,
adversary interleavings included). The cryptographic primitives are
symbolic (perfect); the adversary holds \emph{at most two} of the three
key elements (all subsets explored), can replay any past
container-plus-trap pair, and can falsify either side. Four claims were
verified, each reported with its status and hypotheses:

\begin{itemize}
\item \textbf{P1 (secret below threshold).} With $\le 2$ of the three
elements, the adversary never reaches the learned-secret state:
\emph{proved in the model} (0 counterexamples). The KEK derivation
binds all three elements (HKDF over $K_{user}\,\|\,S_{lockure}\,\|\,
S_{piege}$ with the epoch as context).
\item \textbf{P2 (rollback refusal).} No transition ever decreases the
epoch; replaying a past container-plus-trap pair is refused by the
crossed-epoch check and the highest-seen mirror: \emph{proved in the
model} (structural monotonicity, mechanically re-checked over every
transition).
\item \textbf{P3 (no destruction on wrong passphrase).} Wrong-passphrase
attempts lead only to the benign retry branch (increasing delay); no
reachable state combines wrong-phrase attempts with local destruction:
\emph{proved in the model} (0 counterexamples).
\item \textbf{P4 (blast radius bounded to one epoch).} Proven
falsification destroys only the local container and draft (with a
report); the trap media is the \emph{witness}, never the accused, and
the CF2 fragments coexist: \emph{proved in the model} (0 violations).
\end{itemize}

The limits are named: the checker is an explicit-state explorer
(Tamarin/ProVerif absent from the laboratory --- the deviation is
consigned); this v0.12.0 account was bounded to three epochs, and
the v0.15.0 programme (Section~\ref{sec:v150}) has since lifted the
bound by an inductive invariant, upgrading all four statuses to
\emph{proved unbounded}; side channels, implementation flaws (RAM
discipline, GC),
three-element compromise and quantum adversaries are \emph{out of
model} (hypotheses A1--A5 in the report). ``Proved in the model'' means
exactly that.

\paragraph{Threat model restructured in two planes.} The single threat
table of the specification is restructured along the protocol's own
founding distinction --- obfuscation protects intellectual property,
hardening protects execution --- into two planes:

\begin{itemize}
\item \textbf{Plane IP (intellectual property).} Reverse engineering
(time bought, eroding with LLM decompilers), expropriation (the
evolution of stolen code is blocked --- signed layer-I intent, sealed
inheritance, offline keys, Shamir dead-man's switch --- its execution
is not), redistribution (per-build fingerprints, canary trap).
Confidentiality of the card dominates; loss of confidentiality is
remediable (rotation, variants).
\item \textbf{Plane execution safety.} Card poisoning/corruption
(signature, cross-lock, AEAD), evil-maid on container or trap media
(epoch cross-checks, AEAD tags, the undervivable \code{verif} field),
container-plus-trap rollback (three epoch witnesses), altered binary
(per-block checksums, signed manifest at load), non-conforming
environment (profile checks, fail-closed), agent failure (R2/R5,
oracle from I, budgets). Integrity and availability dominate;
destruction is survivable (CF2 fragments, the trap media).
\end{itemize}

The \emph{intra-Security order} left implicit by the project hierarchy
(D19: Security $>$ Performance $>$ Optimization $>$ Token cost) is now
explicit as project policy: within Security, the IP plane ranks
confidentiality first (remediable losses), while the execution plane
ranks integrity and availability first (irremediable losses); when the
two planes collide, the execution plane wins --- a broken vault protects
nothing. The D81 tension (the control passphrase stored in clear for
survival) and the D183 counter-proposal (encrypted escrow, Shamir
5-of-3) are consigned with the quantified loss/exposure analysis of the
v0.10.0 edition: the expected annual loss under D81 is
$\approx 4.0 \times 10^{-2}\,C_{exposure}$ against
$\approx 1.0 \times 10^{-3}$ under the escrow --- a factor $\approx 40$
--- and D81 \emph{remains} the active decision (the ceremony belongs to
the developer).

\paragraph{The D81 consequence, stated in the model (v0.13.0).} The
honest accounting of the current decision adds a row to the threat
model: the control passphrase lives in clear at the repository (the
survival-over-confidentiality arbitration, D81), and the CF3 container
header carries $S_{serrure}$ --- the machine half of the key material ---
in clear by design. Two of the three key-derivation elements are
therefore \emph{trivially available} to anyone with repository access;
under the threshold adversary of the formal model two elements yield
nothing, but the third --- the trapped medium --- is the only real
confidentiality barrier left, and the integrity guarantees (epoch
witnesses, destruction on proven falsification) carry the protective
load. Two gestures answer this row: the v0.13.0 edition \emph{prepares}
the D183 encrypted-escrow ceremony kit (Shamir 5-of-3 split,
cold-reconstitution proof, media-separation plan, withdrawal checklist)
at the laboratory, \emph{without applying it} --- closing the passphrase
exposure belongs to the developer; and the laboratory's own tests run on
temporary clones with test phrases only (the D105 visible-test boundary:
the real container is never opened by laboratory tests, and test states
are visibly marked as such).

\paragraph{Crystallisation: the statistical model --- corrected (v0.13.0).}
The rule (ten consecutive byte-identical rebuilds before sealing) seals
a build whose per-rebuild divergence probability is $q$ with probability
$(1-q)^n$ (independent rebuilds, the standing hypothesis). Its
operational characteristics: at $q = 50\%$ --- a badly
non-deterministic build --- $n = 5$ would still seal with probability
$3.1\%$ and the historical $n = 10$ with $0.098\%$, a $32\times$
strictness increase, the equivalent of moving from $1.86\sigma$ to
$3.10\sigma$ one-sided; at $q = 1\%$ --- a healthy build --- the rule
seals with probability $95.1\%$ ($n{=}5$) and $90.4\%$ ($n{=}10$). The
values ``$\alpha = 3.29\times10^{-2}$'' and ``$\alpha = 1.08\times
10^{-3}$'' reported in the v0.11.0/v0.12.0 editions were likelihood
ratios $(0.5/0.99)^n$ mislabelled as error probabilities, and their
``sigma'' conversions ($\sqrt{-2\ln p}$) are replaced here by the exact
one-sided normal quantiles. The threshold \emph{remains} ten (project
policy, never changed silently); the corrected analysis is consigned as
the justification and the documented path to any future relaxation.

\paragraph{The replication protocol.} N=1 remains the declared limit;
this edition ships the protocol that would lift it: three independent
axes (a second developer running one full milestone; a second real
application in a different domain; a port to a second non-primary
profile), a T/X coding grid for every defect (five terrain classes T1--T5
--- loader/image, hardware ABI, loader semantics, system conventions,
environment --- against three laboratory classes X1--X3), double
annotation by two independent coders, Cohen's $\kappa$ with the
acceptance threshold $\kappa \ge 0.61$ (substantial agreement), and
pre-registered acceptance criteria. The protocol is proposed, not
executed: the arbitration belongs to the developer.

% ------------------------------------------------------------------
\paragraph{What the campaign adds.}
The D204 field campaign extends the same account in three directions.
First, the locus of failure moved: across 22 deliveries of both arms,
the sealed notation never failed --- the six counted escapes were all
in the rig that delivers and runs it, and every one of them reproduced
a class the study had already named (the unbound import leaf, the
rig-side constructor, the codec). Second, the cost question inverted
on the single instrumented unit: the compiled arm consumed
$\times$1.27 the tokens of the notation arm --- one pair, read from
logs, decomposed by category (the corrective spiral dominates both
arms; the C source costs more tokens than the hand notation), but in
the direction opposite to the registered expectation;
and the verification-time metric, reported for the first time, reads
$23.9\times$ --- H5 falsified on the same single unit, and still
falsified after the present edition's control-preserving optimization
($\times$17.9 measured on the laboratory machine, the critical path
being the emulated contract cryptography itself). Third, the
campaign closed both regimes at once: the controlled comparison
(behaviourally equivalent at the contracts, per the v0.14.0
laboratory result, now field-run at both arms) and the realistic
build (a second real application born under the protocol at both
arms, closed green and signed). The two-way reading stays honest: the
campaign is still one developer, one agent, one machine, and both
metric pairs are $n{=}1$.

\section{Threats to validity}\label{sec:threats}

\emph{Internal.} The oracle is circular with respect to layer I
(both derive from the same specifications); a specification error passes
the oracle by construction --- the v0.11.0 independence runs (distinct
derivation routes, system generator, spec$\leftrightarrow$\allowbreak harness
$\leftrightarrow$\allowbreak hashlib tripartite confrontation) found zero
divergences, and the v0.14.0 programme closed the remaining agent gap
at same-family level: an isolated second agent instance regenerated the
MD5/SHA-256 tables from the mathematical definitions alone with zero
divergences across the 128 constants; true \emph{cross-family}
independence remains the named gap. The crystallisation rule (ten
identical rebuilds) checks
determinism, not correctness. The mutation-testing survivors (56/832
after the v0.14.0 composition re-run; 138/832 at v0.11.0,
concentrated in orchestration/UI blocks without an isolated-mutant
semantic killer) are a named, measured gap --- and the v0.13.0 audit of
our own v0.11.0 benchmark, which caught a perimeter artifact that
inflated a claim fifteen-fold, is itself evidence of this threat class:
instrument defects contaminate results symmetrically, and were caught
here only because the edition required the audit. \emph{External.} N=1: one
project, one developer, one agent family; the compiled-reference
comparison of v0.14.0 exists but shares the laboratory harness and a
single seed; independent replication (a second developer, a second
domain) remains future work. \emph{Construct.} The
escaped-defect count depends on what counts as a ``named'' class --- we
counted conservatively (a class is counted once even if it killed several
runs); the token budgets of M1 and v0.14.0 measure \emph{content reads}
(the latter on the final Loupe-era card, the reasoning estimated and
named), the v0.17.0 re-measurement covers the full instrumented
round-trip on the same card, and the two perimeters are never compared
without the note --- cross-perimeter ratios are meaningless by
construction.
\emph{Construct, campaign-specific: the shared kit and the asymmetric
formats.} The D204 campaign shares its terrain kit between arms while
delivering different formats (an EXE image at a fixed slot for T, a DLL
bound by the Windows loader for C), so the delivery metric mixes
arm-neutral rig defects with format-specific ones. The sensitivity
analysis, per escape: of the six counted escapes, three are
arm-neutral rig classes (a harness constructor, the launcher prose,
the console codec --- a standard flow's test rig can carry exactly
these), and three are manifestations of the single unbound-import-leaf
class, which exists only on the DLL path (an EXE without that import
cannot leave the leaf unbound) --- a compiled-flow-specific escape,
not a notation-versus-compiler difference. The empirical
attribution test (re-delivering the same C piece through a minimal
standard kit) is \emph{prepared but not run}: the extension plan
(D210, pre-registered before any further measurement) fixes it ---
$n \ge 6$ fully instrumented units, identical formats across arms,
randomized arm order, plan hash consigned before the first
measurement. Until then, the two-sided reading stands: zero escapes on
either arm's native code, all six on the rig that delivers and loads
it, three of them format-specific.
\emph{Conclusion.} The terrain is one real Windows machine and one
developer's usage of the application; showcase liveness (0 error-reporter
events) is a weak liveness signal, not a correctness proof. The RISC-V terrain
evidence (M7) was produced under a user-mode emulator with a real
kernel; the \code{fence.i} hardware instruction-cache semantics are
unvalidated on silicon. Memory safety
is \emph{boundedly} verified (2{,}098 accesses classified, zero
stack-smash, 463 named unresolved --- linear analysis, no alias
analysis): a bounded result, not a proof. Constant-time is verified
\emph{structurally} (secret-independent instruction traces); real-hardware
wall-clock behaviour, cache-timing and power side channels remain
UNMEASURED and are marked as not proven --- the v0.15.0 dudect kit is prepared
and waiting for the developer's run, and its laboratory autotest
proves the mechanism only, never the machine. The unbounded proof of
the vault properties rests on an inductive argument whose
epoch-uniformity is structural (the rules compare the epoch only to
1 and take its successor) and machine-witnessed by translations up
to $\delta = 1000$; the checker remains a house-built explorer, not
an independent theorem prover --- the Tamarin/ProVerif absence is
consigned as the deviation.

% ------------------------------------------------------------------
\section{Related work}\label{sec:related}

\paragraph{Positioning (requalified in the present edition).}
The novelty of TRIADE is \emph{not} the generation of assembly by an
LLM --- assembly generation and compiler-style search have been
quantified before: STOKE~\cite{stoke} formulates loop-free binary
superoptimization as stochastic search; SuperCoder~\cite{supercoder}
is the strongest published LLM result on \emph{improving already
compiled} assembly (95.0\% test-pass correctness and a $1.46\times$
average speedup over gcc~-O3 after task-specific fine-tuning ---
superoptimization, a distinct problem from governance);
the LaaC/CompilerEval line~\cite{laac} measures end-to-end compilation
from source code to assembly by LLMs (feasible but with low success
rates --- the compiler's job, attempted whole);
and Nova~\cite{nova} builds generative language models \emph{for}
assembly code (hierarchical attention, contrastive learning, strong
decompilation and similarity results). What no
prior account quantifies is the \emph{governance}: signed-intent
layer-I contracts sealed by a human hand, append-only journals
(engineering, task, decision) as the sole data source, every escaped
defect named to a class and a mechanism before any fix, and a
pre-registered two-arm field campaign comparing the disciplined and
the standard workflows on the same contracts at a real machine ---
with the oracle/terrain epistemology that a laboratory harness is
structurally blind to the classes that only loading, linking and
hardware reveal.

\paragraph{Design by contract and specification-first development.}
TRIADE's layer I applies design-by-contract ideas~\cite{meyer1992} at the
machine-code level: one specification per block, pre/post-conditions,
named failure cases. The novelty is not the contract style but the
binding of contracts to \emph{byte ranges} through the signed card, and
the rule that no notation byte exists before its specification.

\paragraph{N-version and property-based testing.} The differential
oracles and scenario oracles follow the property-based testing
tradition~\cite{claessen2000} and the classic redundancy arguments of
N-version programming~\cite{avizienis1985}: independent instruments (three
concurrent harnesses, oracles regenerated from I, never exported) must
agree. TRIADE extends the independence requirement to the \emph{terrain}
(a machine the laboratory never controls) and to the \emph{rig itself}
(kit defects are named with the same discipline as artifact defects).

\paragraph{Certification discipline.} The milestone/exit-criteria
structure and the append-only, amend-or-invalidate decision log borrow
from airworthiness process standards such as
DO-178C~\cite{do178c}: objectives with independence, traceability, and
configuration control --- here applied to agent-written native code, with
every criterion measured rather than attested.

\paragraph{Supply-chain integrity.} Per-build fingerprints, signed
manifests, sealed tags and key ceremonies align TRIADE with
reproducible builds \cite{lamb2022}, in-toto's layout
attestations~\cite{intoto2019}, SLSA's build levels~\cite{slsa},
survivable-key-compromise update systems (TUF)~\cite{tuf2010} and
automated signing at scale (Sigstore)~\cite{sigstore2022}; the offline
root-key ceremony follows the discipline of DNSSEC key
ceremonies~\cite{rfc4033}. TRIADE's addition is applying these controls
to \emph{fragments of a development state} (the card) rather than to
released artifacts only, and making the sealed tag a fact of the card
while the artifact's bytes stay frozen. Secret sharing of the master
key~\cite{shamir1979} is used both for escrow and for the dead-man's
switch.

\paragraph{Live patching.} Hot patching of native code has a long
lineage, from dynamic software updating~\cite{hicks2002} and
general-purpose DSU systems such as Kitsune~\cite{hayden2012kitsune} and
POLUS~\cite{polusDSU} to industrial kpatch-style
trampolines~\cite{kpatch}. TRIADE's contribution is the
governance wrapper: a patch is a surgical, diffed, card-tracked revision
of blocks at stable addresses, with crystallisation, rollback and
cross-ISA rules (x86-64 trampoline vs.\ RISC-V scratch+\code{fence.i}),
validated by oracle replay and, since M9, requiring a terrain run
recouped on pieces before any seal.

\paragraph{Repository-scale context for LLM agents.} The token-budget
programme (fragmented card, one fragment per operation, no full-binary
reads) addresses the measured context limits of LLM agents: performance
degradation in long contexts~\cite{liu2024} and the need for
repository-scale retrieval and navigation~\cite{zhang2023, jimenez2024}.
The card is, in effect, a \emph{signed repo-map for machine code}: a
compact index that lets an agent operate on a large binary without ever
loading it.

\paragraph{Learned and verified compilation.} The compiler side of the
argument has moved in two directions: \emph{learned} optimisation, with
LLM-based compiler foundation models trained on optimisation
traces~\cite{cummins2024llmcompiler}, and \emph{verified} compilation,
where the translator is not removed but proven, as in
CompCert~\cite{leroycompcert}. TRIADE takes the opposite door: instead
of making the compiler smarter or proving it correct end to end, it
removes the optimising translator from the agent's path and re-derives
its guarantees explicitly --- the WYSIWYG round-trip at every build,
at a measured token cost. The comparison is not adversarial: CompCert
answers ``can the translation be trusted?'', TRIADE answers ``what does
the agent learn when it must write, and therefore own, the bytes?''

\paragraph{Agent-driven software engineering.} SWE-agent~\cite{yang2024sweagent}
quantifies what interface design does for repository-level agent
performance --- at source-code level, in hosted repositories. The RQ
data of the present study is the sealed-notation counterpart: what
breaks when the same class of agent operates on native code under
signed-intent governance. The two layers are complementary; to our
knowledge no quantified account exists between them.

\paragraph{Simulation-to-reality transfer.} The oracle/terrain
epistemology is a deliberate instance of the reality-gap tradition:
evolutionary robotics learned that controllers evolved in simulation
fail on real hardware unless the gap is itself modelled~\cite{jakobi1995reality};
modern sim-to-real work randomises the simulator to close the
gap~\cite{peng2018simtoreal}, and hardware-in-the-loop validation
interposes real components between the two~\cite{bacic2005hil}. TRIADE's
answer is asymmetric: it does not attempt to close the gap by improving
the emulator, but makes the gap itself the measured object --- every
escape is a named, falsified class with an archived evidence piece
(RQ1), and the RQ3 taxonomy is the resulting map of the gap.

\paragraph{Positioning.} None of the cited primitives is novel. The
claimed contribution is the integration --- contract-first notation at
the byte level, signed-intent governance, oracle/terrain epistemology
with named-defect accounting --- evaluated end-to-end by building a real
application with an LLM agent and reporting where it failed.

% ------------------------------------------------------------------
\section{Conclusion}\label{sec:conclusion}

We reported the first academic edition of TRIADE, a protocol under which
an LLM agent writes native code by hand under signed-intent governance,
and a longitudinal single-case evaluation of its eleven-milestone
validation programme. The data, drawn entirely from the protocol's own
append-only records, quantifies the cost (medians of 1{,}277/3{,}374/
3{,}674 tokens per locate/edit/add at the instrumented-round-trip
perimeter; the content-only and M1 values remain labelled history), the
escape dynamics (eleven named defect classes over twelve field
iterations, all environmental; the closing campaign's 22 deliveries
adding zero notation-side escapes and six kit-side ones, every one
reconciled and attributed), and the patch safety (zero
post-patch regression, armed anti-rollback) of the approach. The Loupe
application --- a real Windows hex editor born entirely from the method,
sealed twice, and alive on its user's machine --- is the existence proof
that the pipeline can carry a real deliverable. The second edition strengthened the
evidence: 83.4\% mutation kill with the cryptographic families at
100\%, zero-crash coverage-guided fuzzing of the untrusted-input blocks,
bounded memory-safety verification with zero out-of-frame accesses,
secret-independent instruction traces on the crypto families, an audited
instruction-count parity with a compiled SHA-256 reference (within ten
percent of gcc -O2 and ahead of -Os --- the v0.11.0 ``15.3$\times$''
claim withdrawn by the v0.13.0 audit as a measurement artifact), and zero divergences in oracle-independent regeneration (now including an
isolated-instance regeneration and a binary-search re-derivation). The v0.14.0 empirical programme
adds: a pre-registered compiler-reference comparison showing the sealed
notation observationally equivalent to compiled C from the same contracts
(zero divergence), at instruction parity on the cryptographic core, at a
premium on simple pure blocks, and testable to the same degree; RQ2
re-measured end-to-end with distributions (the v0.17.0 instrumented
round-trip on the native tokenizer, the v0.14.0 content-only values
kept as labelled history under the mandatory perimeter note; the naive
baseline 235--82$\times$ above the card operations, two realistic
baselines quantified alongside); and the mutation campaign re-run with
composition killers (93.3\%, tls13 and sys at 100\%, the 56 remaining
survivors named). The v0.15.0 programme closes the formal side: the four vault
properties proved \emph{unbounded} by an inductive invariant, now
completed by a coverage lemma (every reachable state within the 288
shapes, SMT-discharged with negative controls), a rule-by-rule TLA+
alignment, and a generative conformance campaign of 10{,}000 traces
(one named divergence class, the model's session-falsification gap,
shrunk and analysed; bit-identical at replay), every claim carrying a
systematic status. What remains open
is now precisely measured: the 56 mutation survivors in orchestration
blocks, the single-run scope of the measured hardware constant-time
result (the dudect kit ran twice, signed; more machines and samples
belong to the extension), the unmeasured real-hardware
timing that remains (the \code{fence.i} kit targets a physical RISC-V
board; a third-machine Loupe run is likewise prepared), the missing
independent replication
(the T/X grid and the $\kappa$ calculator are ready; the second
annotator is yet to be recruited) --- and, on the
governance side, the D183 escrow ceremony, prepared and waiting for the
developer's gesture.

% ------------------------------------------------------------------
\section*{Declarations}

\paragraph{Ethics and dual-use.} The artifacts developed under TRIADE
(Loupe, the proxy, the supervisor) are the project's own tools, tested on
the developer's machines; the protocol includes no stealth or
EDR-evasion mechanism, and its specification places furtivity explicitly
out of scope. One counter-measure class is documented honestly in both
documents (the v0.13.0 arbitration, D190): a \emph{light, assumed
anti-debug check by elapsed-time measurement} --- abnormally slow
invocations are treated as dynamic instrumentation --- whose reach is
deterrence, not resistance, and which is consigned as an assumed residual
in the specification's threat table. The work is a dual-use-respectable
engineering study: the same discipline that hardens an artifact against
tampering also makes it auditable. The repository is published under the
Business Source License 1.1 (non-production use until the change date,
then GPL~v3).

\paragraph{Model and toolchain.} The agent is \code{glm-4-plus}
(Z.ai) --- the identifier reported by the API in every instrumented
call since the campaign plan D204; the serving alias is the finest
identifier the API discloses (the exact checkpoint is not exposed:
named, not invented). Sampling parameters are the CLI defaults (not
logged per call: named; the counter logs them going forward).
Verification chain: gcc 14.2.0 (cross mingw-w64), Unicorn 2.1.4,
Capstone 5.0.7, Keystone 0.9.2, Tectonic (XeTeX), TLC/Apalache (the
model-checker fingerprints consigned at the repository).

\paragraph{Licence.} Business Source License 1.1 (Licensor: Alexis
Lopez; Change Date 2030-10-28, converting to GPLv3+). The Additional
Use Grant, ``None'' in the earlier editions (a legal blocker for peer
review and replication), is now \emph{written and applied} at the
repository (D209): the grant permits academic peer review, artifact
evaluation, scientific replication and non-commercial research use of
the Licensed Work, effective 17 November 2026, consigned append-only
in the decision log.

\paragraph{Data availability.} All primary data is public in the project
repository (append-only journal, decision log, evidence vault indexes,
extraction script, generated tables). A versioned archival snapshot
(Zenodo DOI) will be minted at submission time and referenced here ---
the deposit's metadata file is prepared with the present edition
(title, authors, the two archives, SHA-256 fingerprints); the DOI is
inserted upon minting and none is invented before. The
two submission archives of this edition (\code{TRIADE\_soumission\_v0.18.0.zip}
and \code{TRIADE\_complet\_v0.18.0.zip}) contain respectively the
manuscript with its data tables, and the complete sources with raw
extracts and fingerprints; every artifact of the claims tables was
re-verified byte for byte against the regenerated archives (zero
divergence).

\paragraph{Use of AI in this work.} An LLM agent (Z.ai GLM family)
executed the laboratory work of the project under the developer's
instruction, and drafted editorial material under the developer's
direction; all results reported were produced by the protocol's own
instrumentation, and the primary sources (journal, decision log,
evidence vault) were consigned session by session as the work happened.
The agent is not an author. The developer takes full responsibility for
the content of this manuscript.

\paragraph{Acknowledgements.} The author thanks the Z.ai GLM agent
(``the laboratory'') for eleven milestones of disciplined execution, and
the TRIADE decision log for keeping everyone honest.

% ------------------------------------------------------------------
\bibliography{refs}

\end{document}
